\documentclass{../template/note}

\author{Oliver West}
\title{Analysis I}
\date{}

\begin{document}

    \maketitle

    \tableofcontents

    \lecture{Lecture 1 -- 22/01/26}

    Analysis is about things that are small, things that are large, and quantifyig how small or large they are.

    \section{Numerical sequences and series}

    \subsection{Basics}

    \begin{definition}
        A sequence \((x_n)_{n \in \NN}\) on a set \(X\) is an enumerated list where each element is in \(X\).
        
        In this chapter, \(X \subseteq \RR\) or \(X \subseteq \CC\).
    \end{definition}

    What does it mean for a sequence \((x_n)\) to converge to a point \(x \in X\)? What does it mean for \((x_n)\) to diverge to infinity?

    For convergence, we want \(|x_n - x|\) to be small. But what is small? In fact, we want \(|x_n - x|\) to be arbitrarily small, i.e. smaller than any \(\varepsilon > 0\) we choose. What really matters is the \emph{tail} of the sequence -- we can always ignore the first \(N\) terms, and we should pick \(N\) depending on \(\varepsilon\).

    For divergence to infinity, we want \(|x_n|\) to get arbitrarily large, i.e. larger than any \(L > 0\) we choose. Again, we ignore the first \(N\) terms depending on \(L\).
    
    \begin{definition}
        We say \((x_n)\) converges to some finite \(x\), \(x_n \to x\), if, for any \(\varepsilon > 0\), there exists \(N(\varepsilon)\) such that \(|x_n - x| < \varepsilon\) for all \(n \geq N\).

        We say \((x_n)\) diverges to infinity if, for any \(L > 0\), there exists \(N(L)\) such that \(x_n > L\) for all \(n \geq N\), analogously for negative infinity. For complex sequences, we must consider \(|x_n|\).
    \end{definition}
    
    \begin{example}
        \begin{itemize}
            \item Take \(x_n = \frac{1}{n}\), then \(x_n \to 0\) because \(|x_n - 0| = \frac{1}{n} < \varepsilon\) for any \(n \geq N = 1 + \lceil \frac{1}{\varepsilon} \rceil\).
            \item Take \(x_n = \frac{1}{2^n}\), then \(x_n \to 0\) because \(|x_n - 0| = \frac{1}{2^n} < \varepsilon\) for any \(n \geq N = \max(1 + \lceil \log_2(\frac{1}{\varepsilon}) \rceil, 1)\).
            \item Take \(x_n = in\), then \(x_n \to \infty\), since \(\frac{x_n}{i} = n > L\) for any \(n \geq N = 1 + \lceil L \rceil\).
            \item Take \(x_n = (-1)^n\), then \(x_n\) doesn't converge or diverge to infinity.
        \end{itemize}
    \end{example}
    
    \begin{lemma}
        If it exists, the limit of \((x_n)\) is unique.
    \end{lemma}
    
    \begin{proof}
        Suppose \(x_n \to a\) and \(x_n \to b\). Take \(\varepsilon > 0\), then we have
        \begin{align*}
            &\exists N_1 = N_1(\varepsilon) \text{ s.t. } |x_n - a| < \varepsilon \; \forall n \geq N_1 \\
            &\exists N_2 = N_2(\varepsilon) \text{ s.t. } |x_n - b| < \varepsilon \; \forall n \geq N_2 \\
        \end{align*}
        Now take \(N = \max(N_1, N_2)\). We have
        \begin{align*}
            |a-b| &= |a-x_n+x_n-b| \\
            &\leq |x_n-a|+|x_n-b| \\
            &< 2\varepsilon
        \end{align*}
        for all \(\varepsilon > 0\). It follows that \(|a-b| = 0\), so \(a=b\).
    \end{proof}
    
    This justifies the notation \(\lim_{n \to \infty} x_n = x\).

    \begin{proposition}[Sandwich theorems]
        Let \((x_n), (y_n), (z_n)\) be real sequences. Then
        \begin{itemize}
            \item If \(x_n \leq y\) for all \(n\) and \(x_n \to x\), then \(x \leq y\). Note that \(x_n < y\) does not imply \(x < y\).
            \item If \(x_n \to x\), \(z_n \to x\) and \(x_n \leq y_n \leq z_n\) for all \(n\), then \(y_n \to x\).
        \end{itemize}
    \end{proposition}
    
    \begin{proof}
        Left as an exercise.
    \end{proof}
    
    \begin{lemma}
        Let \((x_n)\) be a complex sequence. Then \(x_n \to x\) if and only if \(\Re(x_n) \to \Re(x)\) and \(\Im(x_n) \to \Im(x)\).
    \end{lemma}
    
    \begin{proof}
        Recall that, for \(z \in \CC\), \(|z| = \sqrt{\Re(z)^2 + \Im(z)^2}\).
        \begin{itemize}
            \item[\(\implies\)] Use \(|\Re(z)| \leq |z|\) and \(|\Im(z)| \leq |z|\) with \(z = x_n - x\).
            \item[\(\impliedby\)] Use \(|z| \leq |\Re(z)| + |\Im(z)|\), again with \(z = x_n - x\). Fixing \(\varepsilon > 0\), we can get \(|x_n - x| < 2\varepsilon\), so we're done.
        \end{itemize}
    \end{proof}
    
    \begin{lemma}
        Let \(x_n \to x\), \(y_n \to y\). Then \(x_n + y_n \to x + y\), \(x_ny_n \to xy\).
    \end{lemma}
    
    \begin{proof}
        The first part is analogous to the above result on complex sequences and is an exercise.

        For the second part, we have
        \begin{align*}
            |x_ny_n - xy| &= |x_ny_n - xy_n + xy_n - xy| \\
            &\leq |y_n||x_n-x| + |x||y_n-y|
        \end{align*}
        Now, for \(\varepsilon > 0\), we can pick an \(N\) to make \(|x_n - x| < \varepsilon\), \(|y_n - y| < \varepsilon\). Hence
        \begin{align*}
            |x_ny_n - xy| &< \varepsilon(|x| + |y_n|) \\
            &< \varepsilon(|x| + Y) && \text{by \cref{convergent-bounded}, } Y>0
        \end{align*}
        and so picking \(\varepsilon = \frac{\varepsilon'}{|x| + Y}\) shows the result.
    \end{proof}
    
    \begin{definition}[Bounded sequences]
        We say \((x_n)\) is \emph{bounded} if there exists \(M > 0\) such that \(|x_n| \leq M\) for all \(n\), or, equivalently, if \(\sup |x_n| \leq M\).
    \end{definition}
    
    \begin{lemma}
        \label{convergent-bounded}
        Convergent sequences are bounded.
    \end{lemma}
    
    \begin{proof}
        Take \(\varepsilon = 1\), then there exists \(N\) such that \(|x_n - x| \leq 1\) for any \(n \geq N\), and so \(|x_n| \leq 1 + |x|\) for any \(n \geq N\). But for \(n < N\), \(|x_n| \leq \max_{1 \leq i < N}|x_n| < \infty\).
    \end{proof}
    
    \begin{definition}[Monotone sequences]
        We say a real sequence \((x_n)\) is \emph{monotone} if either
        \begin{itemize}
            \item it is increasing;
            \item or it is decreasing.
        \end{itemize}
    \end{definition}
    
    \begin{proposition}[Monotone convergence theorem]
        \label{monotone-convergence}
        Every bounded monotone sequence converges.
    \end{proposition}
    
    \begin{proof}[Proof sketch]
        Without loss of generality, suppose the sequence is strictly increasing. Using the least upper bound property, we can show that \(x_n \to \sup\{x_n\}\).
    \end{proof}
    
    Can we strengthen to non-monotonic sequences? Clearly not fully, since \(x_n = (-1)^n\) is bounded but not convergent. But, if we delete all the even terms, then it indeed becomes convergent.
    
    \subsection{Bolzano-Weierstrass theorem}

    \begin{theorem}[Bolzano-Weierstrass]
        \label{bw}
        If \((x_n)\) is a real bounded sequence, then it has at least one convergent subsequence.
    \end{theorem}
    
    \begin{definition}[Subsequences]
        Given a sequence \((x_n)\), a subsequence of \((x_n)\) is a sequence of the form \((x_{n_k})_{k \in \NN}\), where \((n_k)\) is a strictly increasing sequence of natural numbers.
    \end{definition}
    
    \begin{lemma}
        If \(x_n \to x\), then any subsequence \((x_{n_k})\) of \((x_n)\) must converge to \(x\).
    \end{lemma}
    
    \begin{proof}
        Since \(n_k < n_{k+1}\), \(n_{k+1} \geq n_k + 1\). By induction, we can show \(n_k \geq k\).

        Now, take \(\varepsilon > 0\), so we have \(N(\varepsilon)\) such that \(|x_n-x| < \varepsilon\) for any \(n \geq N\). So, if \(k \geq N\), then \(n_k \geq k \geq N\), and so \(|x_{n_k}-x|<\varepsilon\), so \(\lim_{k \to \infty}x_{n_k} = x\) as required.
    \end{proof}
    
    \begin{proposition}[Nested interval property]
        \label{nested-interval}
        Let \((I_n)_{n \in \NN}\) be a sequence of closed intervals in \(\RR\) satisfiyng \(I_{n+1} \subseteq I_n\) for every \(n \in \NN\), and suppose \(I_n = [a_n, b_n]\). If \(b_n - a_n = |I_n| \to 0\) as \(n \to \infty\), then \(\bigcap_{n \in \NN} I_n\) contains exactly one point.
    \end{proposition}
    
    \begin{proof}
        We can see that \((a_n)\) and \((b_n)\) are bounded monotone sequences, so we probably want to apply \cref{monotone-convergence}. Indeed, we have \(a_n \to a\), \(b_n \to b\). Since limits preserve inequalities and \(a_n \leq b_n\), we have \(a \leq b\). Now we are ready to show the result.
        \begin{itemize}
            \item \emph{Existence}. For all \(k \geq n\), \(a_k \in [a_k, b_k] \subseteq [a_n, b_n]\), and so \(a_n \leq a_k \leq b_n\). Taking \(k\) to infinity, we have \(a_n \leq a \leq b_n\) for every \(n\), so \(a \in I_n\) for every \(n\), so it is in their intersection.
            \item \emph{Uniqueness}. We know \(b_n - a_n \to b - a\). But we know \(b_n - a_n \to 0\), so \(a = b\). Hence, if \(x\) is in the intersection, \(a_n \leq x \leq b_n\) for all \(n\), so \(a \leq x \leq a\), and hence \(x=a\) as required.
        \end{itemize}
    \end{proof}
    
    \begin{proof}[Proof of theorem]
        We are given \((x_n)\) and \(M > 0\) such that \(|x_n| \leq M\) for any \(n\). We will construct a series of nested intervals from which we can sample our subsequence and apply \cref{nested-interval}.

        Let \(a_1 = -M\), \(b_1 = M\), so \(I_1 = [-M, M] \supset (x_n)\). Now take \(c = \frac{a_1+b_1}{2}\). Then, at least one of the halves \([a_1, c]\) and \([c, b_1]\) will have infinitely many elements of \((x_n)\) -- we take \(I_2\) to be a half that has infinitely many elements. Continue inductively to get \((I_n)\), which are nested by construction and all of which have infinitely many elements of the sequence -- and by \cref{nested-interval}, there exists exactly one \(x\) in their intersection.
        
        We choose \((x_{n_k})\) as follows: pick \(n_1\) such that \(x_{n_1} \in I_1\). Then \(I_2\) has infinitely many elements of \((x_n)\) with index greater than \(n_1\), so pick \(n_2 > n_1\) such that \(x_{n_2} \in I_2\). Continue inductively.

        We have \(x_{n_k} \in I_k\) for every \(k\), so \(x_{n_k} \in \bigcap_{n \leq k} I_n\), and so \(x_{n_k} \to x\).
    \end{proof}

    \begin{corollary}
        The Bolzano-Weierstrass theorem holds in \(\CC\).
    \end{corollary}
    
    
    
    \subsection{Cauchy sequences}

    A Cauchy sequence is a sequence in which the elements are getting closer as the sequence progresses.

    \begin{definition}
        A sequence \((x_n)\) on \(X \subseteq \CC\) is \emph{Cauchy} if, for any \(\varepsilon > 0\), there exists \(N\) such that, for all \(n, m \geq N\), \(|x_n - x_m| < \varepsilon\).
    \end{definition}
    
    \begin{example*}
        \begin{itemize}
            \item Take \(x_n = \frac{1}{n}\). Assume wlog \(m \geq n\). Then
            \begin{align*}
                \left|\frac 1 n - \frac 1 m\right| = \frac 1 n - \frac 1 m = \frac{m-n}{mn} \leq \frac{1}{n} < \eps
            \end{align*}
            and this holds for any \(m \geq n \geq N = \left\lceil\frac1 \eps\right\rceil + 1\). So \((x_n)\) is Cauchy.
            \item Take \(x_n = (-1)^n\). This is not a Cauchy sequence -- take \(n = 2k\), \(m = 2k+1\) for any \(k \in \NN\), so \(|x_n - x_m| = 2\). Hence we can't take \(\eps < 2\).
            \item Take \((x_n)\) on \(\QQ\) defined by truncation of the decimal expansion of \(\sqrt 2\). Assume wlog \(m \geq n\), so
            \begin{align*}
                \left|x_n - x_m\right| < 10^{-(n-1)} \to 0
            \end{align*}
            Note that this sequence does not converge in \(\QQ\).
        \end{itemize}
    \end{example*}
    
    \begin{lemma}
        If \((x_n)\) is Cauchy, then it is bounded.
    \end{lemma}
    
    \begin{proof}
        Take \(\varepsilon = 1\). Then we have \(N\) such that
        \begin{align*}
            |x_n - x_N| < 1
        \end{align*}
        for any \(n \geq N\) (we have fixed \(m = N\)). So \(|x_n| < 1 + |x_N|\).
    \end{proof}
    
    \begin{lemma}
        A complex sequence \((x_n)\) is Cauchy if and only if \((\Re x_n)\) and \((\Im x_n)\) are both Cauchy.
    \end{lemma}    
    
    \begin{lemma}
        If \(x_n \to x\), then \((x_n)\) is Cauchy.
    \end{lemma}
    
    \begin{proof}
        For any \(\eps > 0\), we have \(N\) such that
        \begin{align*}
            |x_n - x_m| &\leq |x_n - x| + |x_m - x| \\
            &\leq 2\varepsilon
        \end{align*}
        for \(n\), \(m \geq N\). This completes the proof.        
    \end{proof}
    Is the converse true? Not in \(\QQ\) at least -- see \(\sqrt 2\) example. But we can still aspire to the converse being true in \(\RR\).

    \begin{theorem}[Completeness of \(\RR\) and \(\CC\)]
        Every Cauchy sequence in \(\RR\) or \(\CC\) converges.
    \end{theorem}
    
    \begin{proof}
        Suppose \((x_n)\) is Cauchy. Then it is bounded, and so by Bolzano-Weierstrass it has a convergent subsequence \((x_{n_k})\) converging to some \(x\). We claim \(x_n \to x\). Observe
        \begin{align*}
            |x_n - x| \leq |x_n - x_{n_k}| + |x_{n_k} - x|
        \end{align*}
        Taking \(\eps > 0\), Cauchy gives \(|x_n - x_{n_k}| < \eps\) for \(n\), \(n_k \geq N_1\), and \(x_{n_k} \to x\) gives \(|x_{n_k} - x| < \eps\) for \(k \geq N_2\). Fixing \(k \geq N_2\) such that \(n_k \geq N_1\), we have \(|x_n - x| < 2\eps\) for every \(n \geq N_1\), which suffices for the result.        
    \end{proof}
    
    \begin{remark*}
        This is nice -- it's our first way to prove convergence without having to guess the limit.
    \end{remark*}
    
    \subsection{Series and convergent series}

    \begin{definition}
        Let \((a_n)\) be a sequence in \(\RR\) or \(\CC\). We say that \(\sum_{n=1}^\infty a_n\) is a \emph{series}, and we say it converges if the sequence of partial sums \(s_k = \sum_{n=1}^k a_n\) converges to some finite \(s\). In this case, \(s\) is called the sum of the series, and we write
        \begin{align*}
            s = \sum_{n=1}^\infty a_n
        \end{align*}        
    \end{definition}
    
    \begin{example*}
        \begin{itemize}
            \item \(\sum_{n=1}^\infty n\) does not converge: \(s_k = \frac{n(n+1)}2\).
            \item geometric series: \(\sum_{n=1}^\infty r^n\) converges iff \(|r| < 1\). We have \(s_k = \frac{1-r^k}{1-r}\) (unless \(r=1\)).
            \item \(\sum_{n=1}^infty \frac{1}{n(n+1)}\) conerges to 1. We have
            \begin{align*}
                s_k = \sum_{n=1}^k \frac{1}{n(n+1)} = \sum_{n=1}^k \left(\frac 1 n - \frac 1 {n+1}\right) = 1 - \frac{1}{k+1} \to 1
            \end{align*}
        \end{itemize}
    \end{example*}
    
    \begin{lemma}
        Fix \(\lambda \in \CC\). If \(\sum a_n\) and \(\sum b_n\) converge, then \(\sum(\lambda a_n + b_n)\) also converges
    \end{lemma}
    
    \begin{proof}
        Exercise.
    \end{proof}
    
    As usual, we should understand that the tail of the series is what is significant to its convergence.

    \begin{lemma}
        If \(a_n = b_n\) for all \(n \geq N\), then \(\sum a_n\) converges if and only if \(\sum b_n\) converges.
    \end{lemma}
    
    \begin{proof}
        We have
        \begin{align*}
            r_k &= \sum_{n=1}^k b_n \\
            &= \sum_{n=1}^{N-1} b_n + \sum_{n=N}^k a_n \\
            &= \underbrace{\sum_{n=1}^{N-1} a_n}_{s_k} + \sum_{n=1}^{N-1} (b_n-a_n)
        \end{align*}
        Taking limits, we conclude that
        \begin{align*}
            \lim_{k \to \infty} r_k = \lim_{k \to \infty} s_k + C
        \end{align*}
        which concludes.        
    \end{proof}
    
    \begin{proposition}[\(n^\text{th}\) term test]
        If \(\sum a_n\) converges, then \(a_n \to 0\).
    \end{proposition}
    
    \begin{remark}
        The converse does not hold. Consider the harmonic series. We have
        \begin{align*}
            s_k = \sum_{n=1}^k \frac 1 n
        \end{align*}
        Observe
        \begin{align*}
            s_{2k} &= s_k + \frac{1}{n+1} + \dots + \frac{1}{n+n} \\
            &\geq s_k + \frac 1 {2n} + \dots + \frac 1 {2n} \\
            &= s_k + \frac{1}{2}
        \end{align*}
        Hence \(s_k\) is not Cauchy, so it does not converge.
    \end{remark}
    
    \begin{proof}
        \(\sum a_n\) converges \(\implies\) \(s_k\) converges, so \((s_k)\) is Cauchy. Then \(s_{k+1} - s_k = a_{k+1} \to 0\).
    \end{proof}
    
    Let's first focus on tests for the convergence of \(\sum a_n\), where \(a_n \geq 0\) for all \(n\).

    \begin{proposition}[Comparison test]
        Suppose we have sequences satisfying \(0 \leq b_n \leq a_n\) for all \(n\). Then \(\sum a_n\) converges \(\implies \sum b_n\) converges.
    \end{proposition}
    
    \begin{proof}
        Set \(s_k\) and \(r_k\) to be the partial sums of \(a_n\), \(b_n\) respectively. Since \(a_n\), \(b_n \geq 0\), \((s_k)\) and \((r_k)\) are increasing. Furthermore, \(s_k \to s \implies s_k \leq s\). But      
        \begin{align*}
            b_n \leq a_n \implies \sum_{n=1}^k b_n \leq \sum_{n=1}^k a_n
        \end{align*}
        so \(r_k \leq s_k\), and hence \(r_k\) converges by the monotone convergence theorem.
    \end{proof}
    
    \begin{example}
        \(\sum_{n=1}^\infty \frac 1 {n^2}\) converges. We have
        \begin{align*}
            \sum_{n=1}^\infty \frac 1 {n^2} = 1 + \sum_{n=2}^\infty \frac 1 {n^2} = 1 + \sum_{n=1}^\infty \frac 1 {(n+1)^2}
        \end{align*}
        But then
        \begin{align*}
            0 \leq \frac 1 {(n+1)^2} \leq \frac 1 {n(n+1)}
        \end{align*}
        and we know \(\sum \frac 1 {n(n+1)} \to 1\). So \(\frac 1 {n^2}\) converges by the comparison test.
    \end{example}
    
    \begin{proposition}[Root test]
        Suppose \(a_n \geq 0\) for all \(n\). Suppose there exists \(a\) such that \(a = \lim_{n \to \infty} a_n^{1/n}\). If
        \begin{itemize}
            \item \(a < 1\), then \(\sum a_n\) converges;
            \item \(a > 1\), then \(\sum a_n\) diverges;
            \item \(a = 1\), then the test is inconclusive.
        \end{itemize}
    \end{proposition}
    
    \begin{proof}
        If \(a > 1\), then we have \(N\) such that \(a_n^{1/n} > 1\) for any \(n \geq N\). But then \(a_n > 1\), so \(a_n \not\to 0\), and we fail the \(n^\text{th}\) term test. Hence the series diverges.

        If \(a < 1\), then we can find \(r \in \RR\) with \(a < r < 1\). Then we have \(N\) such that \(a_n^{1/n} < r\) for any \(n \geq N\). But then \(a_n < r^n\) for all \(n \geq N\). By the comparison test on the tails with a geometric series, \(\sum a_n\) converges.
    \end{proof}
    
    \begin{example}
        \begin{itemize}
            \item Consider \(\sum \frac 1{2^n}\), then \(a_n^{1/n} = \frac{1}{2}\), so \(\sum a_n\) converges.
            \item Consider \(\sum 4^n\), then \(a_n^{1/n} = 4\), so \(\sum a_n\) does not converge.
        \end{itemize}
    \end{example}

    \begin{proposition}[Ratio test]
        If \(a_n \geq 0\) for all \(n\) and \(\frac{a_{n+1}}{a_n} \to a\) for some finite \(a\), then if
        \begin{itemize}
            \item \(a < 1\), then \(\sum a_n\) converges;
            \item \(a > 1\), then \(\sum a_n\) diverges;
            \item \(a = 1\), then the test is inconclusive.
        \end{itemize}
    \end{proposition}
    
    \begin{example}
        \begin{itemize}
            \item If we check \(\sum \frac 1 n\) and \(\sum \frac 1 {n^2}\), we see both lie in the inconclusive regions for both tests.
            \item Let us consider \(\sum \frac{n}{2^n}\). The ratio test gives
            \begin{align*}
                \frac{a_{n+1}}{a_n} = \frac{\frac{n+1}{2^{n+1}}}{\frac{n}{2^n}} = \frac{n+1}{2n} \to \frac 1 2
            \end{align*}
            which is nice. If instead we use the root test, then
            \begin{align*}
                \left(\frac n {2^n}\right)^{\frac 1 n} = \frac{n^{1/n}}{2} \to \frac 1 2
            \end{align*}
            But actually showing that \(n^{1/n} \to 1\) is non-trivial... we need to apply l'H\^opital. Observe
            \begin{align*}
                n^{1/n} = e^{\log(n^{1/n})} = e^{\frac 1 n \log n} \to e^0 = 1
            \end{align*}
            
            
            \item If the ratio test is conclusive for a sequence, then the root test is also conclusive, but not the converse:
            \begin{align*}
                a_n = \begin{cases}
                    2^{-n} & n \text{ even} \\
                    2^{-n+1} & n \text{ odd}
                \end{cases}
            \end{align*}
            Then the limit of ratios does not exist, but the root test shows convergence.            
        \end{itemize}
    \end{example}
    
    \lecture{Lecture \(n\)}

    \begin{proposition}[Integral test]
        Suppose \(f: [1, \infty) \to [0, \infty)\) is a continous decreasing function (and hence integrable in \([1, N]\) for each \(N \in \NN\)), and define \(a_n = f(n)\). Then
        \begin{align*}
            \sum a_n \text{ converges } \iff \lim_{n \to \infty} \int_1^n f(t) \; dt \text{ exists}
        \end{align*}
        Furthermore, as \(k \to \infty\),
        \begin{align*}
            \sum^k a_n - \int_1^k f(x) \; dx \to l \text{ for some }l \in [0, f(1)]
        \end{align*}        
    \end{proposition}

    \begin{proof}
        Observe (DIAGRAM)
        \begin{align*}
            s_k - a_1 &= \sum_{n=2}^k a_n \leq \int_1^k f(t) \; dt \\
            s_{k-1} &= \sum_{n=1}^{k-1} a_n \geq \int_1^k f(t) \; dt
        \end{align*}
        Now
        \begin{itemize}
            \item[\(\implies\)] \(\int_1^k f(t) \; dt \leq s_{k-1} \leq s\) since \(s_k\) is increasing and converges by assumption. Hence \((\int_1^k f(t) \; dt)_k\) converges by MCT.
            \item[\(\impliedby\)] \(s_k \leq a_1 + \int_1^k f(t) \; dt \leq a_1 + \lim_{k \to \infty}\int_1^k f(t) \; dt\). Hence \(s_k\) converges by MCT. 
        \end{itemize}
        For the last part, let
        \begin{align*}
            b_k = \sum^k a_n - \int_1^k f(x) \; dx
        \end{align*}
        We have
        \begin{align*}
            b_k - b_{k-1} = a_k - \int_{k-1}^k f(t) \; dt = f(k) \cdot 1 - \int_{k-1}^k f(t) \; dt \leq 0
        \end{align*}
        and
        \begin{align*}
            \sum_{n=1}^{k-1} a_n &\geq \int_1^k f(t) \; dt \\
            \implies \sum_{n=1}^k - \int_1^k f(t) \; dt &\geq a_k \\
            \implies b_k &\geq a_k = f(k) \geq 0
        \end{align*}
        Hence \(b_k\) is decreasing and bounded below, so converges to some \(l \in \RR\). Since \(0 \leq b_k \leq b_1 = a_1 = f(1)\), we conclude \(0 \leq l \leq f(1)\).
    \end{proof}

    \begin{example}
        \begin{itemize}
            \item \(\sum \frac 1 {n^p}\) converges if and only if \(p > 1\), since
            \begin{align*}
                \lim_{x \to \infty}\int_1^x \frac{1}{t^p} \; dt = \lim_{x \to \infty} \begin{cases}
                    \frac{1}{1-p}x^{1-p} & p \neq 1 \\
                    \log x & p = 1
                \end{cases}
            \end{align*}
            \item \(\sum \frac 1 {n \log n}\) diverges, since
            \begin{align*}
                \int \frac 1 {t \log t} \; dt = \log \log t + C
            \end{align*}
            \item \(\sum \frac 1 {n \log^2 n}\) converges, since
            \begin{align*}
                \int \frac 1 {t \log^2 t} \; dt = -\frac 1 {\log t} + C
            \end{align*}
        \end{itemize}
    \end{example}
    
    \begin{proposition}[Cauchy condensation test, non-examinable]
        Let \(a_n \geq 0\) and \((a_n)\) be decreasing. Then \(\sum a_n\) converges if and only if \(\sum 2^n a_{2^n}\) converges.
    \end{proposition}
    
    \begin{proof}
        Observe
        \begin{align*}
            \sum_{m=1}^{2^n-1} a_n &= a_1 + a_2 + a_3 + \dots + a_{2^n-1} \\
            &\leq a_1 + 2a_2 + \dots + 2^{n-1}a_{2^{n-1}} = \sum_{k=0}^{n-1} 2^k a_{2^k}
        \end{align*}
        and
        \begin{align*}
            \frac{1}{2}\sum_{k=0}^{n-1} 2^k a_{2^k} &= \frac 1 2 a_1 + a_2 + 2a_4 + 4a_8 + \dots + 2^{n-1}a_{2^{n-1}} \\
            &\leq a_1 + a_2 + (a_3 + a_4) + (a_5 + \dots + a_8) + \dots + a_{2^{n-1}} = \sum_{m=1}^{2^{n-1}} a_n
        \end{align*}
        from which the result follows.   
    \end{proof}
    
    
    \begin{proposition}[Alternating series test]
        Let \((a_n)\) be a decreasing sequence, and suppose \(a_n \geq 0\), \(a_n \to 0\). Then \(\sum (-1)^{n+1}a_n\) converges.
    \end{proposition}
    
    \lecture{Lecture \(n+1\)}

    \begin{proof}
        Set \(s_k = \sum^k (-1)^{n+1} a_n\). Then
        \begin{align*}
            s_{2k} = \underbrace{(a_1 - a_2)}_{\geq 0} + \underbrace{(a_3 - a_4)}_{\geq 0} + \dots + \underbrace{(a_{2k-1} - a_{2k})}_{\geq 0}
        \end{align*}
        and so \(s_2 \leq s_4 \leq \cdots\). Likewise, \(s_1 \geq s_3 \geq \cdots\). Moreover, \(s_{2k+1} - s_{2k} = a_{2k+1} \geq 0\), so that
        \begin{align*}
            s_2 \leq s_{2k} \leq s_{2k+1} \leq s_1
        \end{align*}
        so \(s_{2k+1}\) is bounded below by \(s_2\), and \(s_{2k}\) is bounded above by \(s_1\). Hence, by the monotone convergence theorem, we have \(s_{2k} \to s\), \(s_{2k+1} \to \tilde s\). Now, note \(s_{2k+1}-s_{2k} = a_{2k+1} \to 0\), and so \(s = \tilde s\).

        Finally, let us show \(s_n\) converges. We have for every \(\varepsilon > 0\) an \(N_1\) such that \(|s_{2k}-s| < \varepsilon\) for any \(k \geq N_1\), and an \(N_2\) such that \(|s_{2k+1}-s| < \varepsilon\) for any \(k \geq N_2\). But then, taking \(N = 2\max(N_1, N_2)+1\), we can see that \(|s_k - s| < \varepsilon\) for any \(k \geq N\).
    \end{proof}
    
    \begin{proposition}[Dirichlet test]
        Let \((a_n)\) be a decreasing sequence, and suppose \(a_n \geq 0\), \(a_n \to 0\). Let \((b_n)\) be another sequence such that \(s_k = (\sum^k b_n)\) is bounded. Then \(\sum a_nb_n\) converges.
    \end{proposition}
    
    \begin{proof}
        An exercise.
    \end{proof}
    
    \begin{definition}[Absolute convergence]
        We say \(\sum a_n\) is \emph{absolutely convergent} if \(\sum |a_n|\) is convergent.
    \end{definition}
    
    \begin{lemma}
        If \(\sum a_n\) converges absolutely, then \(\sum a_n\) converges.
    \end{lemma}
    
    \begin{proof}
        Let \(s_k = \sum^k a_n\), \(r_k = \sum^k |a_n|\). Observe
        \begin{align*}
            |s_k - s_l| &= \left|\sum_{n=l+1}^k a_n\right| \\
            &\leq \sum_{n=l+1}^k |a_n| \\
            &= r_k - r_l \\
            &= |r_k - r_l|
        \end{align*}
        Hence \(r_k\) convergent \(\implies r_k\) Cauchy \(\implies s_k\) Cauchy \(\implies s_k\) convergent.        
    \end{proof}
    
    \begin{remark*}
        This is awesome, because we can show absolute convergence using all our tests for series with positive terms (the converse is clearly false).
    \end{remark*}
    
    We sometimes call sequences that are convergent but not absolutely convergent conditionally convergent. Conditionally convergent sequences are weird.

    \begin{example}
        Consider
        \begin{align*}
            \sum_{n=1}^\infty \frac{(-1)^{n+1}}{n} = 1 - \frac 1 2 + \frac 1 3 - \frac 1 4 + \frac 1 5 - \frac 1 6 + \cdots
        \end{align*}
        Rearranging terms on the right, we produce
        \begin{align*}
            &\;\left(1 - \frac 1 2\right) - \frac 1 4 + \left(\frac 1 3 - \frac 1 6\right) - \frac 1 8 + \left(\frac 1 5 - \frac 1 {10}\right) - \frac 1 {12} + \left(\frac 1 7 - \frac 1 {14}\right) + \cdots \\
            =&\; \frac 1 2 - \frac 1 4 + \frac 1 6 - \frac 1 8 + \frac 1 {10} - \frac 1 {12} + \frac 1 {14} \\
            =&\; \frac 1 2 \left(1 - \frac 1 2 + \frac 1 3 - \frac 1 4 + \frac 1 5 - \frac 1 6 + \cdots\right)
        \end{align*}
        so we've produced a different value by rearranging the terms. For a conditionally convergent series, the order of the terms matters.       
    \end{example}
    
    \begin{proposition}[Rearrangements under absolute convergence]
        Let \(\sigma: \NN \to \NN\) be a bijection. Let \(a'_n = a_{\sigma(n)}\). Then, if \(\sum a_n\) is absolutely convergent, \(\sum a_n' = \sum a_n\).        
    \end{proposition}
    
    \begin{proof}
        Let \(s_k = \sum^k a_n\), and suppose \(s_k \to s\). Then, for any \(\eps\) we have \(N\) such that, for \(k \geq N\),
        \begin{align*}
            |s_k - s| < \eps \text{ and } \sum_{k > N} |a_n| < \eps
        \end{align*}
        Furthermore, there exists \(N_1 \geq N\) such that \(a_1, \dots, a_N\) are included in \(a_1', \dots, a_N'\). Hence, for \(m \geq N_1\), 
        \begin{align*}
            \sum^m a_n' &= \underbrace{\sum^N a_n}_{s_N} + \;(m-N\text{ terms of }(a_n)\text{ with index}\geq N) \\
            \implies \left|\sum^m a_n'-s\right| &\leq \underbrace{|s_N - s|}_{<\eps} + \underbrace{\sum_{k > N}|a_n|}_{<\eps} \\
            &< 2\eps
        \end{align*}
        which was what we wanted.
    \end{proof}
    
    \section{Continuity of functions}

    \subsection{Limits of functions}

    Take \(f: \CC \to \CC\). We want to define what we mean by \(f(z) \to y\) as \(z \to a\), even if \(a \notin X\). A classic example is \(\frac{\sin x}{x}\) at \(a=0\) -- even though 0 is not in the domain, there are points in the domain that are very close to 0, so it feels like we should be able to assign 0 some value.

    \begin{definition}[Accumulation point]
        Let \(X \subseteq \CC\). We say \(a\) is an \emph{accumulation point} for \(X\) if, for any \(\delta > 0\), there exists \(z \in X \setminus \{a\}\) such that \(|z-a| < \delta\). Otherwise, we say \(a\) is \emph{isolated}.
    \end{definition}
    
    \begin{example}
        \begin{itemize}
            \item Take \(X = \{0\} \cup [1, 2]\) -- then 0 is an isolated point, but all points in \([1, 2]\) are accumulation points.
            \item Take \(S = \{|z| = 1\}\) -- these are accumulation points for the set \(D = \{|z| < 1\}\) (also, every point in \(D\) is an accumulation point for \(D\)).
        \end{itemize}
    \end{example}
    
    \begin{lemma}
        Let \(X \subseteq \CC\), \(a \in \CC\). Then \(a\) is an accumulation point for \(X\) if and only if there exists \((x_n)\) in \(X \setminus \{a\}\) such that \(x_n \to a\).
    \end{lemma}
    
    Let's return to thinking about \(f(z) \to y\) as \(z \to a\). We should be able to make sense of this if \(a\) is either in the domain of \(f\), or is an accumulation point for the domain of \(f\).
    
    \begin{definition}
        Let \(f: X \subseteq \CC \to \CC\). Let \(a \in \CC\) an accumulation point of \(X\). Then we say \(f(z) \to y\) as \(z \to a\) if, for any \(\eps > 0\), there exists \(\delta = \delta(\eps)\) such that \(|z-a| < \delta, z \in X \setminus \{a\} \implies |f(z)-f(a)| < \eps\). We may also write
        \begin{align*}
            \lim_{z \to a}f(z) = y
        \end{align*}

        Furthermore, if \(a\) is an accumulation point (but not in the domain), we say \(f(z) \to \infty\) as \(z \to a\) if
        \begin{itemize}
            \item \(f(z)/|f(z)| \to l\) for some \(l \in \CC\);
            \item and, for any \(L > 0\), there exists \(\delta = \delta(L)\) such that \(|z-a| < \delta, z \in X \setminus \{a\} \implies |f(z)| > L\).
        \end{itemize}
    \end{definition}
    
    \begin{remark}
        If \(a\) is an isolated point of \(X\), then we can always find \(\delta\) sufficiently small such that \(|z-a| < \delta, z \in X \iff z = a\). So we could define a limit for isolated points if we really wanted, but it wouldn't be very interesting.
    \end{remark}
    
    \begin{example}
        Take \(f(x) = \frac{\sin x}{x}\) with domain \(\RR \setminus \{0\}\). We claim \(\frac{\sin x}{x} \to 1\) as \(x \to 0\).
        
        Recall the geometric argument that shows \(\cos x < \frac{\sin x}{x} < 1\) for every \(x \in (0, \frac{\pi}2)\). Hence
        \begin{align*}
            \left|\frac{\sin x}x - 1\right| < \left|1-\cos x\right| = 2\sin^2\left(\frac x 2\right) < \frac{x^2}{2}
        \end{align*}
        Hence, for any \(\eps > 0\), choosing \(\delta = \sqrt{2\eps}\) gives the result.
    \end{example}
    
    \begin{lemma}
        Let \(f: X \subseteq \CC \to \CC\), and take \(a \in \CC\) an accumulation point of \(X\). Then
        \begin{itemize}
            \item \(f(z) \to y\) as \(z \to a\) if and only if \(f(z_n) \to y\) for every sequence \((z_n)\) on \(X \setminus \{a\}\) with \(z_n \to a\);
            \item \(f(z)\) diverges as \(z \to a\) if and only if \(f(z_n)\) diverges for every sequence \((z_n)\) on \(X \setminus \{a\}\) with \(z_n \to a\).
        \end{itemize}
    \end{lemma}
    
    \begin{lemma}
        If \(f: X \subseteq \CC \to \CC\) has a limit at \(a \in \CC\), then this limit is unique.
    \end{lemma}

    \begin{proof}
        Suppose \(f(z) \to x\) and \(f(z) \to y\) as \(z \to a\). Then
        \begin{align*}
            |x-y| &\leq |x-f(z)| + |f(z)-y|
        \end{align*}
        Taking the limit as \(z \to a\), we have
        \begin{align*}
            |x-y| &\leq 0 + 0
        \end{align*}
        so \(x=y\) as required.        
    \end{proof}    
    
    \begin{remark*}
        This justifies the notation \(\lim_{z \to a} f(z) = y\).
    \end{remark*}
    
    \begin{lemma}
        Let \(f\), \(g: X \subseteq \CC \to \CC\), and \(a \in \CC\) an accumulation point of \(X\). Suppose \(\lim_{z \to a}f(z) = y\), \(\lim_{z \to a}g(z) = x\). Then
        \begin{align*}
            \lim_{z \to a} f(z) + g(z) &= y + x \\
            \lim_{z \to a} f(z)g(z) &= yx
        \end{align*}
        If \(g(z) \neq 0\) for any \(z \in X\), and \(x \neq 0\), then
        \begin{align*}
            \lim_{z \to a} \frac{f(z)}{g(z)} = \frac{y}{x}
        \end{align*}
    \end{lemma}
    
    \subsection{Continuity of functions}

    \begin{definition}
        Let \(f: X \subseteq \CC \to \CC\). Then \(f\) is continuous at \(a \in X\) if
        \begin{align*}
            f(a) = \lim_{z \to a}f(z)
        \end{align*}
        or, equivalently, for every \(\eps > 0\), there exists \(\delta = \delta(\eps)\) such that \(|z-a| < \delta\), \(z \in X \implies |f(z)-f(a)| < \eps\).        
    \end{definition}

    If \(a\) is an isolated point, \(f\) is continuous at \(a\), since \(|z-a| < \delta \implies z = a\) for sufficiently small \(\delta\). If, on the other hand, \(a\) is an accumulation point, \(f\) is continuous at \(a\) if and only if
    \begin{align*}
        \lim_{z \to a} f(z) = f(a)
    \end{align*}
    

    \begin{definition}
        Let \(f: X \subseteq \CC \to \CC\). We say \(f\) is continuous if it is continuous at every point in \(X\).
    \end{definition}

    \begin{example}
        \begin{itemize}
            \item \(f(z) = z\) is continuous -- take \(\delta = \eps\).
            \item Consider \begin{align*}
                f(x) = \begin{cases}
                    \sin\left(\frac 1 x\right) & x \neq 0 \\
                    0 & x = 0
                \end{cases}
            \end{align*}
            This function is not continuous. Take \(x_n = \frac{1}{(2n+\frac 1 2)\pi}\), which satisfies \(x_n \to 0\) -- but \(f(x_n) = \sin\left(\left(2n+\frac 1 2\right)\pi\right) = 1\) for any \(n\), so \(f(x_n) \to 1 \neq f(0)\).
        \end{itemize}
    \end{example}
    
    \lecture{Lecture 8 -- 10/02/26}
    
    There is another, equivalent notion of continuity which is extremely useful.

    \begin{definition}[Sequential continuity]
        Let \(f: X \subseteq \CC \to \CC\). We say \(f\) is \emph{sequentially continuous} at \(a \in X\) if \(f(z_n) \to f(a)\) for all \((z_n)\) on \(X\) with \(z_n \to a\).
    \end{definition}
    
    \begin{proposition}
        Let \(f: X \subseteq \CC \to \CC\). Then \(f\) is continuous if and only if it is sequentially continuous.
    \end{proposition}
    
    \begin{proof}
        \begin{itemize}
            \item[\(\implies\)] By continuity of \(f\), for any \(\eps > 0\) we have \(\delta\) such that \(|z-a| < \delta\), \(z \in X\) implies \(|f(z) - f(a)| < \eps\). Take \((z_n)\) on \(X\) with \(z_n \to a\), then there exists \(N = N(\delta(\eps))\) such that, for any \(n \geq N\), \(|z_n-a| < \delta\) and hence \(|f(z_n)-f(a)| < \eps\).
            \item[\(\impliedby\)] Let us take the contrapositive. If \(f\) is not continuous at \(a\), then there exists \(\eps > 0\) suchb that, for any \(\delta > 0\), one can find \(z = z(\delta) \in X\) such that \(|f(z)-f(a)| \geq \eps\) while \(|z-a| < \delta\). Now take \(\delta = 1, \frac 1 2, \dots, \frac 1 n, \dots\) to find \((z_n)\) with \(|z_n-a| < \frac 1 n\), so \(z_n \to a\), but \(|f(z_n)-f(a)| \geq \eps\) for all \(n\) and a non-zero \(\eps\). Hence \(f\) cannot be sequentially continuous, as required.
        \end{itemize}
    \end{proof}
    
    \begin{example*}
        Let us have some more examples.
        \begin{itemize}
            \item Consider the Dirichlet function \(f(x) = 1_\QQ(x)\). This is discontinuous at every real \(a\):
            \begin{itemize}
                \item if \(a \in \QQ\), we can find \((x_n) \subset \RR \setminus \QQ\) with \(x_n \to a\), but of course \(f(x_n) = 0\), so \(f(x_n) \to 0 \neq 1 = f(a)\);
                \item if \(a \in \RR \setminus \QQ\), we proceed analogously.
            \end{itemize}
            \item \(f(x) = \sin x\) is continuous as a real function. Fix \(a \in \RR\) and, for some \(\eps > 0\), pick \(\delta = \min(\eps, \pi)\). Now
            \begin{align*}
                |f(x)-f(a)| &= |\sin x - \sin a| \\
                &= 2\left|\cos\left(\frac{x+a}2\right)\sin\left(\frac{x-a}2\right)\right| \\
                &\leq 2\left|\sin\left(\frac{x-a}2\right)\right| \\
                &\leq |x-a| < \eps
            \end{align*}
            since \(\frac{x-a}{2} \in \left[-\frac \pi 2, \frac \pi 2\right]\).
        \end{itemize}
    \end{example*}
    
    Now let us look at some important properties of continuous functions.    

    \begin{lemma}
        Let \(f\), \(g: X \subseteq \CC \to \CC\) be continuous at \(a \in \CC\). Then, so are \(\lambda f + g\) for \(\lambda \in \CC\) fixed, \(fg\) and, if \(f(z) \neq 0\) for any \(z \in X\), \(1/f(z)\).
    \end{lemma}
    
    \begin{lemma}
        Let \(U\), \(V \subseteq \CC\), and \(f: U \to V\), \(g: V \to \CC\). Suppose \(f\) is continuous at \(a \in U\) and \(g\) is continuous at \(f(a) \in V\) -- then, \(g \circ f: U \to \CC\) is continuous at \(a\).
    \end{lemma}
    
    \begin{proof}
        Since \(g\) is continuous at \(f(a)\), we have \(\sigma = \sigma(\eps)\) such that \(y \in V\), \(|y-f(a)| < \sigma\) implies \(|g(y) - g(f(a))| < \eps\). Furthermore, since \(f\) is continuous at \(a\), we have for any \(\sigma\) a \(\delta = \delta(\sigma) = \delta(\eps)\) such that
        \begin{align*}
            z \in V, |z-a| < \delta \implies |f(z) - f(a)| < \sigma \implies |g(f(z)) - g(f(a))| < \eps
        \end{align*}
    \end{proof}
    

    \subsection{Extreme value theorem}

    \begin{definition}[Closed sets]
        A set \(X \subseteq \CC\) is \emph{closed} if all sequences \((x_n)\) in \(X\) which converge in \(\CC\) have limits in \(X\).
    \end{definition}
    
    \begin{example*}
        \begin{itemize}
            \item The closed interval \([0, 1]\) is closed, while the open interval \((0, 1)\) is not.
            \item \([0, 1)\) and \((0, 1]\) are not closed.
            \item \([0, \infty)\) is closed, while \((0, \infty]\) is not.
        \end{itemize}
    \end{example*}

    \begin{definition}[Bounded sets]
        We say \(X \subseteq \CC\) is \emph{bounded} if there exists \(M > 0\) such that \(X \subseteq \{z \in \CC: |z| \leq M\}\), or in other words \(\sup_{z \in X} |z| \leq M\).
    \end{definition}
    
    \begin{example*}
        \([0, 1]\), \((0, 1)\), \([0, 1)\) and \((0, 1]\) are all bounded, while \([0, \infty)\) and \((0, \infty)\) are not.
    \end{example*}
    
    \begin{proposition}
        \label{cts-preserves-cbdd}
        Let \(X \subseteq \CC\) be a closed bounded set. If \(f: X \to \CC\) is continuous, then \(f(X) \subseteq \CC\) is a closed bounded subset of \(\CC\).
    \end{proposition}

    \begin{proof}
        \begin{itemize}
            \item \emph{\(f(X)\) is bounded}. Suppose not. Then, for each \(n \in \NN\), we can find \(x_n \in X\) such that \(|f(x_n)| > n\). Now, \((x_n)\) is a sequence in \(X\), so it must be bounded, so Bolzano-Weierstrass applies -- there exists a convergent subsequence \((x_{n_k})\) such that \(x_{n_k} \to x\) and, by closedness of \(X\), \(x \in X\). On the other hand, \(|f(x_{n_k})| > n_k \geq k\), so \(f(x_{n_k})\) cannot converge as \(k \to \infty\). But then \(f\) is not continuous, since it is not sequentially continuous.
            \item \emph{\(f(X)\) is closed}. Take \((y_n)\) in \(f(X)\), and suppose it converges to some \(y \in \CC\) -- we want to show \(y \in f(X)\). Since \(y_n \in f(X)\), we have \(x_n \in X\) such that \(f(x_n) = y_n\), and \((x_n)\) has a subsequence \((x_{n_k})\) converging to \(x \in X\) by the closure of \(X\) and the Bolzano-Weierstrass theorem. By continuity of \(f\),
            \begin{align*}
                y = \lim_{k \to \infty} y_{n_k} = \lim_{k \to \infty} f(x_{n_k}) = f\left(\lim_{k \to \infty} x_{n_k}\right) = f(x)
            \end{align*}
            which was what we wanted.            
        \end{itemize}
        
        
    \end{proof}
    

    \begin{theorem}[Extreme value theorem]
        \label{evt}
        Let \(X \subseteq \RR\) be a closed bounded set, and let \(f: X \to \RR\) be continuous. Then there exist \(a\), \(b \in X\) satisfying
        \begin{align*}
            f(a) = \sup f(X); \qquad f(b) = \inf f(X)
        \end{align*}
    \end{theorem}

    \lecture{Lecture 9 -- 12/02/26}
    
    \begin{proof}
        We will consider specifically the supremum, since the infimum is analogous.

        Let \(M = \sup f(X)\), which exists since \(f(X)\) is bounded by \cref{cts-preserves-cbdd}. By definition, \(M - \frac 1 n\) is not an upper bound on \(f(X)\) for any \(n \in \NN\), that is, we can find a sequence \((y_n)\) with elements in \(f(X)\) such that \(M - \frac 1 n < y_n \leq M\), which produces a sequence \((x_n)\) on \(X\) such that \(M - \frac 1 n < f(x_n) \leq M\). Now take limits and use the fact that \(f\) is continuous -- we have \(M \leq \lim_{n \to \infty} f(x_n) \leq M\).

        Finally, by \cref{cts-preserves-cbdd}, we know that, \(f(X)\) is closed, so, since \((y_n)\) is a sequence in \(f(X)\), its limit is also in the image of \(X\).        
    \end{proof}

    \begin{remark*}
        We know that the sup and inf exist -- the key point is that \(f\) actually achieves both of these.
    \end{remark*}
    
    \subsection{Intermediate value theorem}
    
    \begin{theorem}[Intermediate value theorem]
        \label{ivt}
        If \(f: [a, b] \to \RR\) is continuous then, for any \(y\) satisfying \(f(a) \leq y \leq f(b)\), there exists \(x\) satisfying \(a \leq x \leq b\) such that \(f(x) = y\).
    \end{theorem}
    
    \begin{proof}
        If \(f\) is constant, \(y=f(a)\) or \(y=f(b)\), then the result is clear. If not, wlog let \(f(a) < y < f(b)\). Set \(S = \{x \in [a, b]: f(x) \leq y\}\). Then \(a \in S\), so \(S\) is non-empty -- furthermore, \(S\) is bounded above by \(b\), so we have \(c = \sup S \in [a, b]\). We claim that \(f(c) = y\). Observe:
        \begin{itemize}
            \item We claim \(f(c) \leq y\). Suppose not. Then \(\eps = f(c) - y > 0\). By continuity of \(f\), we have \(\delta\) such that \(|x-c| < \delta\), \(x \in [a, b]\) implies \(|f(x) - f(c)| < \eps\), so \(f(x) > f(c) - \eps = y\), so \(x \notin S\). This means \((c - \delta, c) \cap S = \emptyset\), which contradicts the definition of \(c\) (in particular that \(c\) is the least upper bound).
            \item We claim \(f(c) \geq y\). Suppose not. Then \(\eps = y - f(c) > 0\). By continuity of \(f\), we have \(\delta\) such that \(|x-c| < \delta\), \(x \in [a, b]\) implies \(|f(x) - f(c)| < \eps\), so \(f(x) < f(c) + \eps = y\). But then \(f(c + \frac{\delta}2) < y\), so \(c \leq c + \frac{\delta}2 \in S\), which contradicts the definition of \(c\) (in particular that \(c\) is an upper bound).
        \end{itemize}
        The result follows.        
    \end{proof}
    
    \begin{remark*}
        A nice interpretation of this theorem is that continuous functions from reals to reals map intervals to intervals.
    \end{remark*}
    
    \begin{example*}
        We can use this result to show the existence of \(n^\text{th}\) roots. Take \(t > 0\) and \(N \in \NN\), and consider
        \begin{align*}
            f: [0, 1 + t] &\to \RR \\
            x &\mapsto x^N
        \end{align*}
        This is continuous, and \(f(0) = 0 < t < (1+t)^N = f(1+t)\). Hence, there exists a \(c\) satisfying \(c^N = t\) by the intermediate value theorem. Note that we do not have uniqueness.        
    \end{example*}
    
    \begin{definition}[Monotone functions]
        A function \(f: [a, b] \to \RR\) is \emph{monotone} if it is either increasing, \(x_1 \leq x_2 \implies f(x_1) \leq f(x_2)\); or decreasing, \(x_1 \leq x_2 \implies f(x_2) \leq f(x_1)\).
    \end{definition}
    
    \begin{proposition}[Inverse function theorem]
        \label{ift-1}
        Let \(f: [a, b] \to \RR\) be continuous and strictly monotone. Let \(c = \min(f(a), f(b))\), \(d = \max(f(a), f(b))\). Then \(f: [a, b] \to [c, d]\) is a bijection, and \(f^{-1}: [c, d] \to [a, b]\) is continuous and strictly monotone.
    \end{proposition}

    \begin{proof}
        Since \(f\) is continuous, it achieves its extrema by \cref{evt}; since it is monotone, it achieves them at its endpoints, so it maps \([a, b] \to [c, d]\) -- since it is strictly monotone, it does so injectively. By \cref{ivt}, it is surjective, and hence it is bijective.

        For the rest of the proof, we suppose wlog that \(f\) is strictly increasing. Suppose for contradiction that \(f^{-1}: [c, d] \to [a, b]\) is not strictly increasing -- then we have \(y_1, y_2 \in [c, d]\) with \(y_1 < y_2\) and \(f^{-1}(y_1) \geq f^{-1}(y_2)\). But then, applying \(f\) to both sides, we obtain \(y_1 = f(f^{-1}(y_1)) \geq f(f^{-1}(y_2)) = y_2\), which is a contradiction.

        Finally, it remains to show \(f^{-1}\) is continuous. Let \(y_0 \in [c, d]\) and \(x_0 = f^{-1}(y_0)\). We split into three cases.
        \begin{itemize}
            \item \(x_0 \in (a, b)\). Fix \(\eps > 0\), and choose \(\eta \in (0, \eps]\) sufficiently small for \([x_0 - \eta, x_0 + \eta] \subseteq [a, b]\). Now, since \(f\) is strictly increasing,
            \begin{align*}
                x_0 - \eta &< x_0 < x_0 + \eta \\
                \implies f(x_0 - \eta) &< f(x_0) = y_0 < f(x_0 + \eta)
            \end{align*}
            Let us now take \(\delta = \min(f(x_0 + \eta) - y_0, y_0 - f(x_0 - \eta)) > 0\). Suppose \(|y-y_0| < \delta\), then
            \begin{align*}
                f(x_0 - \eta) \leq y_0 - \delta < y &< y_0 + \delta \leq f(x_0 + \eta) \\
                \implies f(x_0 - \eta) < y &< f(x_0 + \eta) \\
                \implies |f^{-1}(y) - x_0| &< \eta \\
                \implies |f^{-1}(y) - f^{-1}(y_0)| &< \eps
            \end{align*}
            as required.
            \item \(x_0 = a\), \(y_0 = f(a)\). Fix \(\eps > 0\), and choose \(\eta = \min(\eps, b-a)\). Set \(\delta = f(a + \eta) - f(a) > 0\). Then \(|y-y_0| < \delta\), \(y \in [c, d]\) gives
            \begin{align*}
                f(a) \leq y &< f(a + \eta) \\
                \implies a \leq f^{-1}(y) &< a + \eta \\
                \implies |f^{-1}(y)-f^{-1}(y_0)| &< \eta \leq \eps
            \end{align*}
            as required.
            \item \(x_0 = b\) proceeds analogously.               
        \end{itemize}
        This completes the result.        
    \end{proof}
    
    
    
    \lecture{Lecture 10 -- 14/02/26}

    \section{Differentiation}

    \subsection{Basics}

    \begin{definition}[Derivative]
        Let \(f: X \subseteq \CC \to \CC\), \(a \in X\). We say \(f\) is differentiable at \(a\) if the limit
        \begin{align*}
            \lim_{x \to a} \frac{f(x)-f(a)}{x-a} = \lim_{h \to 0} \frac{f(a+h)-f(a)}{h}
        \end{align*}
        exists. We call this limit the derivative of \(f\) at \(a\), \(f'(a)\) or \(\frac{df}{dx}(a)\).        
    \end{definition}
    
    \begin{remark*}
        We can't really make sense of the derivative at an isolated point of the domain. For accumulation points, we can distinguish
        \begin{itemize}
            \item interior points: we can approach from every direction and we need the limit to be direction-independent;
            \item non-interior points: the limit is `domain-restricted'.
        \end{itemize}
    \end{remark*}
    
    \begin{example*}
        \begin{itemize}
            \item \(f(z) = z\) is differentiable everywhere:
            \begin{align*}
                f'(z) = \lim_{h \to 0} \frac{f(z+h)-f(z)}{h} = \lim_{h \to 0} 1 = 1
            \end{align*}
            \item \(f(z) = \bar z\) is not differentiable anywhere. Indeed,
            \begin{align*}
                h = \lambda \in \RR: \lim_{\lambda \to 0} \frac{f(z+\lambda)-f(z)}{\lambda} &= \lim_{\lambda \to 0} \frac{\overline{z+\lambda}-\bar z}{\lambda} \\
                &= \lim_{\lambda \to 0} \frac{\lambda}{\lambda} = 1 \\[10pt]
                h = i\lambda, \, \lambda \in \RR: \lim_{\lambda \to 0} \frac{f(z+i\lambda)-f(z)}{i\lambda} &= \lim_{\lambda \to 0} \frac{\overline{z+i\lambda}-\bar z}{i\lambda} \\
                &= \lim_{\lambda \to 0} \frac{-i\lambda}{i\lambda} \\
                &= -1
            \end{align*}
            Hence, the limit is different depending on the direction we approach from, and so does not exist.
            \item \(f(z) = \sin z\) is differentiable. Observe
            \begin{align*}
                f'(x) &=^? \lim_{h \to 0} \frac{\sin(x+h)-\sin x}{h} \\
                &= \lim_{h \to 0} \frac{\sin h \cos x}{h} + \lim_{h \to 0} \frac{\sin x(\cos h-1)}{h} \\ 
                &= cos x \lim_{h \to 0} \frac{\sin h}{h} + \sin x \lim_{h \to 0} \frac{\cos h - 1}{h}
            \end{align*}
            The limit of \(\frac{\sin h}{h}\) is 1 as already established, and the limit of \(\frac{\cos h-1}{h}\) can be shown to be zero by, for example, using a double angle formula.  
        \end{itemize}
    \end{example*}

    \begin{lemma}
        Let \(f\), \(g: X \subseteq \CC \to \CC\) be differentiable at \(a \in X\). Then so are \(f+g\), \(fg\) and, if \(f(z) \neq 0\) for any \(z \in X\), so is \(1/f\). Furthermore, we have
        \begin{align*}
            (f+g)'(a) &= f'(a) + g'(a) \\
            (fg)'(a) &= f'(a)g(a) + f(a)g'(a) \\
            \left(\frac 1 f\right)'(a) &= -\frac{f'(a)}{f(a)^2}
        \end{align*}
    \end{lemma}
    
    \begin{proof}
        The addition rule is an exercise. For the multiplication rule, we have
        \begin{align*}
            (fg)'(a) &= \lim_{h \to 0} \frac{(fg)(a+h) - (fg)(a)}{h} \\
            &= \lim_{h \to 0} \frac{f(a+h)g(a+h)-f(a)g(a)}{h} \\
            &= \lim_{h \to 0} \frac{\left(f(a+h)-f(a)\right)g(a+h) + \left(f(a+h)-g(a)\right)f(a)}{h} \\
            &= \lim_{h \to 0} \frac{\left(f(a+h)-f(a)\right)g(a+h)}{h} + f(a)\lim_{h \to 0} \frac{g(a+h)-g(a)}{h} \\
            &= \lim_{h \to 0} g(a+h) f'(a) + g'(a)f(a)
        \end{align*}
        We will later show that differentiability implies continuity, which will allow us to conclude.
    \end{proof}    

    \begin{example*}
        Let us consider some examples.
        \begin{itemize}
            \item \(f(z) = z^2 = z \cdot z\) is differentiable with \(f'(z) = 2z\). By induction, \(z^n\) is differentiable with \(f'(z) = nz^{n-1}\), so polynomials are differentiable.
            \item \(f(z) = \frac 1 z\) is differentiable on \(\CC \setminus \{0\}\) with \(f'(z) = -\frac 1 {z^2}\). By induction, \(\frac 1 {z^n}\) has derivative \(-\frac n {z^{n+1}}\).
        \end{itemize}
        More generally, rational functions \(\frac{p(x)}{q(x)}\) with \(p\), \(q\) polynomials are differentiable wherever \(q(x) \neq 0\).        
    \end{example*}

    \lecture{Lecture 11 -- 17/02/26}
    
    \begin{proposition}[Chain rule]
        Let \(U\), \(V \subseteq \CC\) and \(f: U \to V\), \(g: V \to \CC\). Suppose that \(f\) is differentiable at \(a \in U\) and \(g\) is differentiable at \(f(a) \in V\). Then \(g \circ f: U \to \CC\) is differentiable at \(a \in U\), and \((g \circ f)'(a) = g'(f(a))f'(a)\).
    \end{proposition}
    
    We will first need a lemma.

    \begin{lemma}
        \label{little-o-derivative}
        Let \(f: X \subseteq \CC \to \CC\). Then \(f\) is differentiable at \(a \in X\) if and only if there exists \(A \in \CC\) and \(\xi: \{z: z + a \in X\} \to \CC\) satisfying \(\xi(h) \to 0\) as \(h \to 0\) such that
        \begin{align*}
            f(a+h) = f(a) + Ah + \xi(h)|h|
        \end{align*}
        Furthermore, \(A = f'(a)\).
    \end{lemma}
    
    \begin{remark*}
        In other words, \(f(a+h) \approx f(a) + Ah\), and the function \(\xi(h)\) allows us to quantify that error. In little o notation, this statement is \(f(a+h) = f(a) + Ah + o(h)\).
    \end{remark*}
    
    \begin{proof}[Proof of lemma]
        \begin{itemize}
            \item[\(\impliedby\)] We have
            \begin{align*}
                \lim_{h \to 0} \frac{f(a+h)-f(a)}h &= \lim_{h \to 0}\frac{Ah + \xi(h) |h|}{h} \\
                &= A + \lim_{h \to 0} \left(\xi(h) \frac{|h|}{h}\right) \\
                &= A
            \end{align*}
            since \(\frac{|h|}{h}\) is bounded and \(\xi(h) \to 0\) as \(h \to 0\). Hence the limit exists and \(f\) is differentiable as required.
            \item[\(\implies\)] Let \(A = f'(a)\), so
            \begin{align*}
                \lim_{h \to 0} \frac{f(a+h)-f(a)}{h} &= A \\
                \implies \lim_{h \to 0} \frac{f(a+h)-f(a)-Ah}{h} &= 0
            \end{align*}
            Now take
            \begin{align*}
                \xi(h) = \begin{cases}
                    \frac{f(a+h)-f(a)-Ah}{|h|} & h \neq 0 \\
                    0 & h = 0
                \end{cases}
            \end{align*}
            which satisfies \(\xi(h) \to 0\) as \(h \to 0\).
        \end{itemize}
    \end{proof}
    
    \begin{proof}[Proof of proposition]
        Since \(f\), \(g\) are differentiable at \(a\), \(f(a)\) respectively, there exist error functions \(\xi_f\), \(\xi_g\) satisfying \(\xi_f(h) \to 0\) as \(h \to 0\) and \(\xi_g(k) \to 0\) as \(k \to 0\), and
        \begin{align*}
            f(a+h) &= f(a) + f'(a)h + \xi_f(h)|h| \\
            g(f(a)+k) &= g(f(a)) + g'(f(a))k + \xi_g(k)|k|
        \end{align*}
        Now
        \begin{align*}
            (g \circ f)(a+h) - (g \circ f)(a) &= g(f(a+h)) - g(f(a)) \\
            &= g(f(a) + \underbrace{f'(a)h + \xi_f(h)|h|}_k) - g(f(a)) \\
            &= g(f(a)) + g'(f(a))\left(f'(a)h + \xi_f(h)|h|\right) + \xi_g(k)|k| - g(f(a)) \\
            &= g'(f(a))f'(a)h + g'(f(a))\xi_f(h)|h| + \xi_g(k)|k|
        \end{align*}
        so we want to make
        \begin{align*}
            \xi(h) &= g'(f(a))\xi_f(h) + \xi_g(k)\frac{|k|}{|h|} \\
            &= g'(f(a))\xi_f(h) + \xi_g(hf'(a)+\xi_f(h)|h|)\cdot|f'(a)+\xi_f(h)| \\
        \end{align*}
        It remains only to check that \(\xi(h) \to 0\) as \(h \to 0\). We have
        \begin{align*}
            \lim_{h \to 0} \xi(h) &= \lim_{h \to 0} \left(g'(f(a))\xi_f(h) + \xi_g(hf'(a)+\xi_f(h)|h|)\cdot|f'(a)+\xi_f(h)|\right) \\
            &= \lim_{h \to 0} g'(f(a))\xi_f(h) + \lim_{h \to 0} \xi_g(hf'(a)+\xi_f(h)|h|) \cdot \lim_{h \to 0} |f'(a)+\xi_f(h)|
        \end{align*}
        Then the first term is zero because \(\xi_f(h) \to 0\) as \(h \to 0\), and the first multiplicand in the second term goes to 0 as \(h \to 0\) because it is clear that its argument goes to 0 as \(h \to 0\).
    \end{proof}
    
    \begin{example*}
        Take
        \begin{align*}
            f(x) = \begin{cases}
                x \sin \left(\frac 1 x\right) & x \neq 0 \\
                0 & x = 0
            \end{cases}
        \end{align*}
        At \(x \neq 0\), the chain and product rules give
        \begin{align*}
            f'(x) &= 1 \cdot \sin\left(\frac 1 x \right) + x \cdot \cos\left(\frac 1 x\right) \cdot \left(-\frac 1 {x^2}\right) \\
            &= \sin \left(\frac 1 x\right) - \frac 1 x \cos\left(\frac 1 x\right)
        \end{align*}
        At \(x = 0\), the definition gives
        \begin{align*}
            f'(x) = \lim_{h \to 0} \frac{f(h)-f(0)}h = \lim_{h \to 0} \frac{f(h)} h = \lim_{h \to 0} \sin \left(\frac 1 x\right)
        \end{align*}
        which does not exist -- hence the function is not differentiable at 0.
    \end{example*}
    
    \begin{lemma}
        \label{diff-cont}
        If \(f: X \subseteq \CC \to \CC\) is differentiable at \(a \in X\), then it is continuous at \(a \in X\).
    \end{lemma}
    
    \begin{proof}
        By \cref{little-o-derivative}, we have
        \begin{align*}
            \lim_{x \to a} f(x) = \lim_{h \to 0} f(a+h) &= \lim_{h \to 0} f(a) + Ah + \xi(h)|h| \\
            &= f(a) + A \lim_{h \to 0} h + \lim_{h \to 0} \left(\xi(h)|h|\right) \\
            &= f(a)
        \end{align*}
    \end{proof}

    \lecture{Lecture 12 -- 19/02/26}
    
    \subsection{Mean value theorems}

    We think of the derivative as an instantaneous rate of change, and we want to try and relate this to \emph{average} rates of change.
    Let us first have a proposition-lemma.
    
    \begin{proposition}[Rolle's theorem]
        \label{rolle}
        Let \(f\) cts\(: [a, b] \to \RR\) be differentiable on \((a, b)\), and suppose \(f(a) = f(b)\). Then there exists \(c \in (a, b)\) such that \(f'(c) = 0\).
    \end{proposition}
    
    \begin{proof}
        By \cref{evt}, \(f\) attains its min and max, that is there exist \(x_m\) and \(x_M \in [a, b]\) with
        \begin{align*}
            \sup_{[a, b]} f = f(x_M) \qquad \qquad \inf_{[a, b]} f = f(x_m)
        \end{align*}
        If \(f\) is constant, then \(f'(c) = 0\) for any \(c \in (a, b)\), so we are done. Otherwise, it either holds that \(f(x_M) > f(a)\) (so \(x_M \in (a, b)\)) or \(f(x_m) < f(a)\) (so \(x_m \in (a, b)\)). It reamins to show that, at either of these points, \(f' = 0\). We take cases.
        \begin{itemize}
            \item \(x_m \in (a,b)\). Then, for sufficiently small \(h\),
            \begin{align*}
                \sign\left(\frac{f(x_m + h)-f(x_m)}h\right) = \sign(h)
            \end{align*}
            so, taking \(h \to 0\) from the right and left, we conclude \(f'(x_m)\) is both positive and negative respectively, i.e. \(f'(x_m) = 0\).
            \item \(x_M \in (a, b)\). Then we have
            \begin{align*}
                \sign\left(\frac{f(x_M + h)-f(x_M)}h\right) = -\sign(h)
            \end{align*}
            which again shows \(f'(x_M) = 0\) by taking limits on both sides.
        \end{itemize}
    \end{proof}
    
    \begin{theorem}[Mean value theorem]
        \label{mvt}
        Let \(f\) cts\(: [a, b] \to \RR\) be differentiable on \((a, b)\). Then, there exists \(c \in (a, b)\) such that
        \begin{align*}
            f'(c) = \frac{f(b)-f(a)}{b-a}
        \end{align*}
    \end{theorem}

    \begin{proof}
        Let
        \begin{align*}
            l(x) = f(a) + \frac{f(b)-f(a)}{b-a}(x-a)
        \end{align*}
        in other words, the secant line connecting \((a, f(a))\) to \((b, f(b))\). Then let \(\phi(x) = f(x) - l(x)\), so \(\phi(a) = \phi(b) = 0\). Hence Rolle gives \(c\) satisfying \(\phi'(c) = 0\), that is
        \begin{align*}
            f'(c) - \frac{f(b)-f(a)}{b-a} &= 0 \\
            \implies f'(c) &= \frac{f(b)-f(a)}{b-a}
        \end{align*}
        as required.
    \end{proof}
    
    \begin{corollary}
        \label{mvt-cor-1}
        Let \(f\) cts\(: [a, b] \to \RR\) be differentiable on \((a, b)\). Given \(h\) such that \(a + h \in [a, b]\), there exists \(\theta \in (0, 1)\) such that
        \begin{align*}
            f(a+h) = f(a) + hf'(a+\theta h)
        \end{align*}
    \end{corollary}
    
    \begin{proof}
        Restrict the domain of \(f\) to \([a, a+h]\). Then, by the mean value theorem, we have \(c \in [a, a+h]\) with
        \begin{align*}
            f'(c) &= \frac{f(a+h)-f(a)}{h} \\
            \implies f(a+h) &= f(a) + hf'(c)
        \end{align*}
        Then, let \(\theta\) be such that \(c = a+\theta h\).        
    \end{proof}
    
    \begin{corollary}
        \label{mvt-cor-2}
        Let \(f\) cts \(: [a, b] \to \RR\) be differentiable on \((a, b)\). Then
        \begin{itemize}
            \item \(f' \geq 0\) on \((a,b) \implies f\) is increasing, resp. strictly increasing;
            \item \(f' \leq 0\) on \((a,b) \implies f\) is decreasing, resp. strictly decreasing;
            \item \(f' = 0\) on \((a,b) \implies f\) is constant.
        \end{itemize}
    \end{corollary}
    
    \begin{proof}
        Exercise.
    \end{proof}
    
    \begin{remark*}
        Statements such as \cref{mvt-cor-2} do not necessarily hold if the domain \(X\) is an arbitrary subset of \(\RR\). For instance, take
        \begin{align*}
            f: \QQ &\to \RR \\
            x &\mapsto \begin{cases}
                0 & x^2 > 2 \\
                1 & x^2 < 2
            \end{cases}
        \end{align*}
        which is continuous (yes, really!) and satisfies \(f'(x) = 0\) for any \(x \in \QQ\), but is not constant.        
    \end{remark*}
    
    \begin{lemma}
        Let \(f: \CC \to \CC\) be differentiable (\(\implies\) continuous) in \(\CC\). Then, \(f'(z) = 0\) for all \(z \in \CC \implies f\) constant.
    \end{lemma}
    
    \begin{proof}
        We want to reduce to the \(\RR\) case. For some \(z \in \CC\), take \(g(t) = f(tz)\) where \(g: [0, 1] \to \CC\). Note \(g\) is continuous on \([0, 1]\), \(g\) is differentiable on \((0, 1)\), so we can apply \cref{mvt} and its corollaries to \(\Re(g(t))\) and \(\Im(g(t))\):
        \begin{align*}
            g'(t) = zf'(tz) = 0 \implies \begin{cases}
                \Re(g'(t)) = 0 \\
                \Im(g'(t)) = 0
            \end{cases}
        \end{align*}
        and so \(g(t)\) is constant -- hence \(f(0) = g(0) = g(1) = f(z)\). But \(z\) was arbitrary, so \(f(z) = f(0)\) for any \(z \in \CC\).
    \end{proof}
    
    \begin{remark*}
        We can use similar ideas to show a version of the mean value theorem in \(\CC\).
    \end{remark*}
    
    \begin{theorem}[Inverse function theorem -- differentiable version]
        \label{ift-2}
        Let \(f\) cts\(: [a, b] \to \RR\) be differentiable on \((a, b)\). Assume that \(f'(x) > 0\) for any \(x \in (a, b)\). Then \(f: [a, b] \to [f(a), f(b)]\) is a bijection with \(f^{-1}\) cts\(: [f(a), f(b)] \to [a, b]\) differentiable on \((f(a), f(b))\) with
        \begin{align*}
            (f^{-1})'(y) = \frac 1 {f'(f^{-1}(y))}
        \end{align*}
    \end{theorem}
    
    \begin{proof}
        By \cref{mvt-cor-2}, \(f\) is strictly increasing, so our previous version of IFT, \cref{ift-1}, gives that \(f\) is a bijection and has a continuous inverse \(f^{-1}\). It remains to show that \(f^{-1}\) is differentiable and satisfies the above.

        Let \(y \in (f(a), f(b))\) and \(x = f^{-1}(y)\). Given \(h\) such that \(y + h \in (f(a), f(b))\), define \(k\) such that \(y+h=f(x+k)\), or equivalently \(k=f^{-1}(y+h)-x\). Then
        \begin{align*}
            \frac{f^{-1}(y+h)-f^{-1}(y)}{h} &= \frac{x+k-x}{f(x+k)-y} \\
            &= \frac{k}{f(x+k)-f(x)} \\
            &= \frac 1{\frac{f(x+k)-f(x)}{k}}
        \end{align*}
        so the result holds so long as \(k \to 0\) as \(h \to 0\) -- then we can take limits on both sides. But
        \begin{align*}
            \lim_{h \to 0} k &= \lim_{h \to 0} f^{-1}(y+h) - f^{-1}(y) \\
            &= 0
        \end{align*}
        because \(f^{-1}\) is continuous.       
    \end{proof}
    
    \begin{example}
        Let us consider some examples.
        \begin{itemize}
            \item Fix \(R > 0\), and define
            \begin{align*}
                f: [0, R] &\to \RR \\
                x &\to x^N
            \end{align*}
            so \(f\) is continuous on \([0, R]\) and differentiable on \((0, R)\) with \(f'(x) = Nx^{N-1}\). Hence, by \cref{ift-2}, \(f: [0, R] \to [0, R^N]\) is a bijection, and \(f^{-1}(y) = y^{1/N}\) is differentiable with
            \begin{align*}
                (f^{-1})'(y) = \frac 1 {f'(f^{-1}(y))} = \frac 1 {N (f^{-1}(y))^{N-1}} = \frac 1 N y^{\frac 1 N - 1}
            \end{align*}
            \item \(f(x) = e^x\) is continuous and differentiable on \(\RR\) with \(f'(x) = e^x > 0\). Hence, we have an inverse \(f^{-1}(y) = \log y\) such that
            \begin{align*}
                (f^{-1})'(y) = \frac 1 {e^{f^{-1}(y)}} = \frac 1 {e^{\log y}} = \frac 1 y
            \end{align*}
        \end{itemize}
    \end{example}
    
    \begin{proposition}[Cauchy's mean value theorem]
        Let \(f, g\) cts\(:[a, b] \to \RR\) be differentiable on \((a, b)\). Then there exists \(c \in (a, b)\) such that
        \begin{align*}
            g'(c) \left(f(b)-f(a)\right) = f'(c)\left(g(b)-g(a)\right)
        \end{align*}
    \end{proposition}
    
    \begin{proof}
        We will reduce to Rolle's theorem, \cref{rolle}. Let
        \begin{align*}
            \phi(x) &= \left(g(x)-g(a)\right)\left(f(b)-f(a)\right) - \left(g(b)-g(a)\right)\left(f(x)-f(a)\right) \\
            \implies \phi'(x) &= g'(x)\left(f(b)-f(a)\right) - f'(x)\left(g(b)-g(a)\right)
        \end{align*}
        and \(\phi(a) = \phi(b) = 0\). The result follows.        
    \end{proof}
    
    \begin{proposition}[L'H\^opital's rule, informal statement]
        For certain limits expressed as a ratio of functions \(f\) and \(g\), one can calculate the limit by taking derivatives first.
    \end{proposition}
    
    \lecture{Lecture 13 -- 21/02/26}

    \subsection{Higher derivatives and Taylor's theorem}
    
    \begin{definition}[Higher derivatives]
        Suppose \(f: X \subseteq \CC \to \CC\) is differentiable on \(X\). We say \(f\) is \emph{\(n\)-times differentiable on \(X\)} if \(f\) is \((n-1)\)-times differentiable on \(X\) and
        \begin{align*}
            f^{(n-1)}: X &\to \CC \\
            x &\mapsto f^{(n-1)}(x)
        \end{align*}
        is differentiable.
    \end{definition}
    
    \begin{definition}[Continuous differentiability]
        We say \(f: X \subseteq \CC \to \CC\) is \(k\)-times \emph{continuously differentiable} on \(X\), and write \(f \in C^k(X)\), if \(f\) is \(k\)-times differentiable on \(X\) and \(f^{(k)}\) is continuous.
    \end{definition}
    
    \begin{remark*}
        Recall that \(f\) differentiable \(\implies f\) continuous, so \(f\) \(k\)-times differentiable \(\implies f\) \((k-1)\)-times continuously differentiable.
    \end{remark*}
    
    \begin{definition}[Smoothness]
        We say \(f: X \subseteq \CC \to \CC\) is \emph{smooth} if it is \(k\)-times differentiable for any \(k\), or, equivalently, if \(f \in C^k(X)\) for any \(k\).
    \end{definition}
    
    Since the derivative of a function at a point \(x_0\) determines the slope of the tangent, \(f(x_0)\) and \(f'(x_0)\) together yield a linear approximation to the graph of \(f\) at \(x_0\):
    \begin{align*}
        f(x) \approx f(x_0) + f'(x_0)(x-x_0)
    \end{align*}
    for \(x \approx x_0\). If we have more derivatives, however, we should in principle be able to get a better approximation:
    \begin{align*}
        f'(x) &\approx f'(x_0) + f''(x_0)(x-x_0) \\
        \implies f(x) &\approx f(x_0) + f'(x_0)(x-x_0) + \frac 1 2 f''(x_0)(x-x_0)^2 
    \end{align*}
    and, if \(f\) is \(n\)-times differentiable,
    \begin{align*}
        f(x) \approx \sum_{k=0}^n \frac 1 {k!} f^{(k)}(x_0)(x-x_0)^k
    \end{align*}
    for \(x \approx x_0\).    
    
    \begin{definition}[Taylor polynomials and Taylor series]
        If \(f: X \subseteq \CC \to \CC\) is \(n\)-times differentiable and \(x_0 \in X\), then
        \begin{align*}
            T_{n,f,x_0}(x) = \sum_{k=0}^n \frac 1 {k!} f^{(k)}(x_0)(x-x_0)^k
        \end{align*}
        is called the \emph{\(n^\text{th}\)-order Taylor polynomial of \(f\) at \(x_0\)}. If \(f\) is smooth, then
        \begin{align*}
            T_{f,x_0}(x) = \sum_{k=0}^\infty \frac 1 {k!} f^{(k)}(x_0)(x-x_0)^k
        \end{align*}
        is called the \emph{Taylor series of \(f\) at \(x_0\)}.        
    \end{definition}
    
    We would like to quantify exactly how good this approximation is.

    \begin{theorem}[Taylor's theorem, Lagrange remainder]
        \label{taylor-lagrange}
        Suppose \(f: X \subseteq \RR \to \RR\) and its \(n-1\) derivatives are continuous on \([a, a+h]\) and \(f\) is \(n\)-times differentiable on \((a, a+h)\). Set
        \begin{align*}
            R_{n, f, a}(h) = f(a+h) - \sum_{k=0}^{n-1} \frac 1 {k!} f^{(k)}(a) h^k = f(a+h) - T_{n-1, f, a}(h)
        \end{align*}
        Then, there exists \(\theta(h) \in (0, 1)\) such that
        \begin{align*}
            R_{n, f, a}(h) = \frac 1 {n!} h^n f^{(n)}(a+\theta h) 
        \end{align*}
    \end{theorem}
    
    \begin{remark*}
        The conditions for the above are a consequence of \(f \in C^n((c, d))\) for any \((c, d) \supset [a, a+h]\). In this case, \(f^{(n)}\) is also bounded on \([a, a+h]\), so there exists \(M_n = \sup_{[a, a+h]}|f^{(n)}|\). Then
        \begin{align*}
            |R_{n,f,a}(h)| \leq M_n \frac{h^n}{n!}
        \end{align*}
    \end{remark*}

    Morally speaking, we are saying that \(R_{n, f, a}(h) = O(h^n)\) as \(h \to 0\), and we know exactly what the constant is. Unfortunately, this does \emph{not} tell us anything about whether \(R_{n, f, a}(h) \to 0\) as \(n \to \infty\): even if \(f \in C^\infty\), we have no idea how \(M_n\) changes with \(n\).

    \begin{remark*}
        We should also note that there is nothing important about \(h > 0\) -- if \(h = -t < -0\), and \(f\) and its \(n-1\) derivatives are continuous and \([a-t, a]\) and \(f\) is \(n\)-times differentiable on \((a-t, a)\), then we can apply the above to \(g(x) = f(-x)\) to obtain
        \begin{align*}
            g(-a+t) &= \sum_{k=0}^{n-1} \frac{(-1)^k}{k!} f^{(k)}(a) t^k + \frac{(-1)^n}{n!}t^n f^{(n)}(a-\theta t) \\
            \implies f(a-t) &= \sum_{k=0}^{n-1} \frac{1}{k!} f^{(k)}(a)(-t)^k + \frac{1}{n!}(-t)^n f^{(n)}(a+\theta(-t))
        \end{align*}
    \end{remark*}
    
    \begin{proof}
        Without loss of generality, we may take \(a=0\) (otherwise apply to \(g(x) = f(x+a)\)). Let
        \begin{align*}
            \phi: [0, h] &\to \RR \\
            t &\mapsto f(t) - \sum_{k=0}^{n-1} \frac{1}{k!} f^{(k)}(0)t^k - \frac{t^n}{n!}B = f(t) - T_{n-1,f,0}(t) - \frac{t^n}{n!}B
        \end{align*}
        where \(B\) is chosen such that \(\phi(h) = 0\). Note \(\phi\) and its first \(n-1\) derivatives are continuous on \([0, h]\), and its \(n^\text{th}\) derivative exists on \((0, h)\). One can see fairly easily that \(\phi^{(k)}(0) = 0\) for \(k = 0, \dots, n-1\). We can thus apply \cref{mvt-cor-1} repeatedly:
        \begin{align*}
            \phi(0) = \phi(h) = 0 &\implies \exists \theta_1 \in (0, 1) \text{ s.t. } \phi'(\theta_1h) = 0 \\
            \phi'(0) = \phi'(\theta_1h) = 0 &\implies \exists \theta_2 \in (0, 1) \text{ s.t. } \phi''(\theta_2\theta_1h) = 0 \\
            \vdots \\
            \phi^{(n-1)}(0) = \phi^{(n-1)}(\theta_{n-1}\cdots\theta_1h) = 0 &\implies \exists \theta_n \in (0, 1) \text{ s.t. } \phi^{(n)}(\theta_n\cdots\theta_1h) = 0 \\
        \end{align*}
        Now, setting \(\theta = \theta_n \cdots \theta_1 \in (0, 1)\), we end up with \(\theta \in (0, 1)\) such that \(\phi^{(n)}(\theta h) = f^{(n)}(\theta h) - B = 0\), so \(B = f^{(n)}(\theta h)\) as required.
    \end{proof}
    
    \begin{proof}[Proof (alternative)]
        Again wlog let us take \(a=0\). Now let
        \begin{align*}
            g: [0, h] &\to \RR \\
            t &\mapsto f(h) - \sum_{k=0}^{n-1} \frac 1 {k!} f^{(k)}(t) (h-t)^k = f(h) - T_{n-1,f,0}(h-t)
        \end{align*}
        Note that \(g\) is continuous on \([0, h]\) and differentiable on \((0, h)\) with
        \begin{align*}
            g'(t) &= -\sum_{k=1}^{n-1} \frac 1 {(k-1)!} f^{(k)}(t) (-1) (h-t)^{k-1} - \sum_{k=0}^{n-1} \frac 1 {k!} f^{(k+1)}(t) (h-t)^k \\
            &= -\sum_{k=0}^{n-2} \frac 1 {k!} f^{(k+1)}(t)(h-t)^{k} - \sum_{k=0}^{n-1} \frac 1 {k!} f^{(k+1)}(t) (h-t)^k \\
            &= -\frac 1 {(n-1)!} f^{(n)}(t) (h-t)^{n-1}
        \end{align*}
        For \(p \in \{1, \dots, n\}\), set
        \begin{align*}
            \phi_p(t) = g(t) - \left(\frac{h-t}{h}\right)^pg(0)
        \end{align*}
        Then \(\phi_p(0) = \phi_p(h) = 0\) for each \(p\), so, for each \(p\), there exists \(\theta \in (0, 1)\) with
        \begin{align*}
            0 &= \phi_p'(\theta h) = g'(\theta h) + \frac{p(1-\theta)^{p-1}}{h}g(0) \\
            \implies g(0) &= -\frac{hg'(\theta h)}{p(1-\theta)^{p-1}}
        \end{align*}
        But we can see \(g(0) = R_{n, f, 0}(h)\), so
        \begin{align*}
            R_{n, f, 0}(h) = -\frac{hg'(\theta h)}{p(1-\theta)^{p-1}} = \frac{h^n}{p(n-1)!}(1-\theta)^{n-p}f^{(n)}(\theta h)
        \end{align*}
        so putting \(p = n\) gives the result.
    \end{proof}
    
    This alternative proof approach will become extremely useful for our next theorem.

    \begin{theorem}[Taylor's theorem, Cauchy remainder]
        \label{taylor-cauchy}
        Let \(f\) and \(R_{n,f,a}\) be as in \cref{taylor-lagrange}. Then there exists \(\theta \in (0, 1)\) such that
        \begin{align*}
            R_{n, f, a}(h) = h^n \frac{(1-\theta)^{n-1}}{(n-1)!} f^{(n)}(a + \theta h)
        \end{align*}
    \end{theorem}

    \begin{proof}
        This follows immediately from taking \(p=1\) in the final line of the above.
    \end{proof}
    
    \lecture{Lecture 14 -- 24/02/26}
    
    \begin{remark*}
        Supposing \(f \in C^n((c, d))\) for some \((c, d) \supset [a, a+h]\), we again have \(R_{n,f,a}(h) = O(h^n)\). However, as we will see, the form of the remainder is often better behaved in \(n\).
    \end{remark*}
    
    \begin{example*}
        Consider \(f(x) = x^q\), \(q \in \QQ\), which is smooth on \((0, \infty)\) with
        \begin{align*}
            f^{(n)}(x) = q (q-1) \dots (q-k+1) x^{q-k}
        \end{align*}
        so that
        \begin{align*}
            T_{n-1,f,1}(x) &= \sum_{k=0}^{n-1} \frac{q(q-1) \dots (q-k+1)}{k!}x^k \\
            &= \sum_{k=0}^{n-1} \binom q k x^k
        \end{align*}
        and we consider
        \begin{align*}
            R_{n,f,1}(x) = (1+x)^q - \sum_{k=0}^{n-1} \binom q k x^k
        \end{align*}
        for \(x\) in the interval \((-1, 1)\). Let us compare both our remainder options:
        \begin{itemize}
            \item \emph{Lagrange remainder}. We have \(\theta \in (0, 1)\) such that
            \begin{align*}
                R_{n,f,1}(x) = \binom q n (1+\theta x)^{q-n}x^n
            \end{align*}
            Then, if \(n \geq q\), \((1+\theta x)^{q-n} = \frac 1{(1+\theta x)^{n-q}} \leq 1\) for \(x \in [0, 1)\), so
            \begin{align*}
                |R_{n,f,1}(x)| \leq \left|\binom q n x^n\right| \to 0
            \end{align*}
            as \(n \to \infty\). But this argument fails for \((-1, 0)\).
            \item \emph{Cauchy remainder}. We have \(\theta \in (0, 1)\) such that
            \begin{align*}
                R_{n,f,1}(x) &= (1-\theta)^{n-1}n\binom q n (1+\theta x)^{q-n}x^n \\
                &= (1-\theta)^{n-1}q\binom {q-1}{n-1} (1+\theta x)^{q-n}x^n \\
                &= q \binom{q-1}{n-1} \left(\frac{1-\theta}{1+\theta x}\right)^{n-1} (1+\theta x)^{q-1}x^n
            \end{align*}
            Now \(1+\theta x = (1-\theta) + \theta(x+1) \geq 1-\theta\), so \(\left(\frac{1-\theta}{1+\theta x}\right)^{n-1} \leq 1\) for \(n \geq 1\). Hence
            \begin{align*}
                |R_{n,f,1}(x) &\leq |qx|\left|\binom{q-1}{n-1}\right||1+\theta x|^{q-1} \\
                &\leq C\left|\binom{q-1}{n-1}\right|
            \end{align*}
            by taking the sup over all \(x \in (-1, 1)\). But this goes to 0 as \(n \to \infty\). So in fact the Cauchy remainder does give us a stronger result.    
        \end{itemize}
    \end{example*}
    
    \begin{definition}[Analyticity, non-examinable]
        We say \(f \in C^\infty((c, d))\) is \emph{analytic} if, for every \(a \in (c, d)\), there exists an \(r > 0\) such that, for any \(|h| < r\), 
        \begin{align*}
            f(a+h) = \lim_{n \to \infty} T_{n, f, a}(h) \iff \lim_{n \to \infty} R_{n,f,a}(h) = 0
        \end{align*}
        in other words, \(f\) is equal to its Taylor series.
    \end{definition}
    
    \begin{example*}
        As we spent the last page showing, \(x^q\) is analytic in \((0, 2) \iff (1+x)^q\) is analytic in \((-1, 1)\).
    \end{example*}
    
    \begin{remark*}
        Note that, even if \(R_{n,f,a}(h)\) does not go to zero, the Taylor series might still converge -- just not to the actual function.
    \end{remark*}
    
    \section{Integration of functions}

    \subsection{Basics}

    We want to define \(\int_a^b f(x) \D x\) as the signed area under a graph, which is easy if it's, say, a rectangle. Otherwise, we can try approximating it with slices.
    
    \begin{definition}[Upper and lower Riemann sums]
        Let \(f: [a, b] \to \RR\) be bounded. Given a partition \(\mathcal P = \{x_0, \dots, x_n\}\) with \(a = x_0 < \dots < x_n = b\) of \([a, b]\), we set
        \begin{align*}
            L(f, \mathcal P) &= \sum_{j=1}^n (x_j-x_{j-1})\inf_{x \in [x_{j-1},x_j]}f(x)\\
            U(f, \mathcal P) &= \sum_{j=1}^n (x_j-x_{j-1})\sup_{x \in [x_{j-1},x_j]}f(x)
        \end{align*}
        to be the \emph{lower} and \emph{upper Riemann sums} respectively.
    \end{definition}
    
    \begin{example*}
        Let
        \begin{align*}
            f: [0, 1] &\to {0, 1} \\
            x &\mapsto \begin{cases}
                1 & x \in \QQ \\
                0 & x \in \RR \setminus \QQ
            \end{cases}
        \end{align*}
        No matter what \(\mathcal P\) we take, each of our intervals will contain both a rational and an irrational, so
        \begin{align*}
            U(f, \mathcal P) &= 1 \\
            L(f, \mathcal P) &= 0
        \end{align*}
    \end{example*}
    
    \begin{remark*}
        For \(f\) bounded, i.e. \(\sup_{[a, b]} |f| = k\) for some \(k\),
        \begin{align*}
            U(f, \mathcal P) &\leq \sum_{j=1}^n (x_j-x_{j-1})\sup_{[a,b]}|f| \\
            = k(b-a)
        \end{align*}
        Analogously,
        \begin{align*}
            L(f, \mathcal P) &\geq \sum_{j=1}^n (x_j-x_{j-1})(-\sup_{[a,b]}|f|) \\
            = -k(b-a)
        \end{align*}
        Also, it's clear that \(U(f, \mathcal P) \geq L(f, \mathcal P)\). Hence
        \begin{align*}
            &\{L(f, \mathcal P): \mathcal P \text{ a partition of }[a,b]\} \\
            &\{U(f, \mathcal P): \mathcal P \text{ a partition of }[a,b]\}
        \end{align*}
        are both bounded sets and (clearly) non-empty.        
    \end{remark*}
    
    \begin{definition}[Upper and lower integrals]
        Let \(f: [a, b] \to \RR\) be bounded. We say
        \begin{align*}
            I_*(f) &= \sup_{\mathcal P} L(f, \mathcal P) \\
            I^*(f) &= \inf_{\mathcal P} U(f, \mathcal P)
        \end{align*}
        are the \emph{lower} and \emph{upper integrals} respectively.        
    \end{definition}
    
    \begin{definition}[Riemann integrability]
        We say \(f\) is \emph{Riemann integrable} if \(I_* = I^*\). Then we say
        \begin{align*}
            I(f) = \int_a^b f(x) \D x = I_*(f) = I^*(f)
        \end{align*}
    \end{definition}
    
    \begin{remark*}
        By convention, \(\int_a^b f = -\int_b^a f\).
    \end{remark*}
    
    \begin{example*}
        Take \(f\) as in the previous example. Then \(I_*(f) = 0\) and \(I^*(f) = 1\), so \(f\) is not Riemann integrable.
    \end{example*}
    
    \lecture{Lecture 15 -- 26/02/26}
    
    The coarsest possible partition for \([a, b]\) is \(\{a, b\}\). Our intuition suggests that finer partitions lead to a better estimate for the integral.

    \begin{lemma}
        Let \(\mc P\) and \(\mc P'\) be partitions of \([a, b]\), and suppose \(\mc P \supseteq \mc P'\), so \(\mc P'\) is a refinement of \(\mc P\). Then 
        \begin{align*}
            L(f, \mc P) \leq L(f, \mc P') \leq U(f, \mc P') \leq U(f, \mc P)
        \end{align*}
        for every bounded \(f: [a, b] \to \RR\).
    \end{lemma}
    
    \begin{proof}
        Let us induct. Take \(\mc P = \{x_0, \dots, x_n\}\) and pick \(y \in (x_{k-1}, x_k)\) for some \(x\) -- then, put \(\mc P' = P \cup \{y\}\). Then
        \begin{align*}
            \sup_{x \in [x_{k-1},y]} f(x), \, \sup_{x \in [y,x_k]} f(x) &\leq \sup_{x \in [x_{k-1}, x_k]}f(x) \\ 
            \implies (y-x_{k-1})\sup_{x \in [x_{k-1},y]} f(x) + (x_k-y)\sup_{x \in [y,x_k]} f(x) &\leq (x_k-x_{k-1}) \sup_{x \in [x_{k-1}, x_k]}f(x)
        \end{align*}
        so that \(U(f, \mc P') \leq U(f, \mc P)\) as required. Similar reasoning with the infimum would suffice for the lower sums. Hence we can induct on the elements of \(\mc P' \setminus \mc P\).
    \end{proof}
    
    \begin{lemma}
        Let \(\mc P\) and \(\mc P'\) be any partitions of \([a, b]\). Then
        \begin{align*}
            L(f, \mc P) \leq U(f, \mc P')
        \end{align*}
        for every bounded \(f: [a, b] \to \RR\).
    \end{lemma}
    
    \begin{proof}
        Let \(\mc P'' = \mc P \cup \mc P'\), so \(\mc P''\) is a refinement of both \(\mc P\) and \(\mc P'\). Then
        \begin{align*}
            L(f, \mc P) \leq U(f, \mc P'') \leq U(f, \mc P')
        \end{align*}
    \end{proof}
    
    \begin{lemma}
        For a bounded \(f: [a, b] \to \RR\),
        \begin{align*}
            I_*(f) \leq I^*(f)
        \end{align*}
    \end{lemma}

    \begin{proof}
        We know that
        \begin{align*}
            L(f, \mc P) &\leq \inf_{\mc P'} U(f, \mc P') \\
            \implies \underbrace{\sup_{\mc P} L(f, \mc P)}_{I_*(f)} &\leq \underbrace{\inf_{\mc P'} U(f, \mc P')}_{I^*(f)}
        \end{align*}
    \end{proof}
    
    \begin{remark*}
        It follows that \(f\) is integrable if and only if \(I_*(f) \geq I^*(f)\).
    \end{remark*}
    
    \begin{proposition}[Riemann integrability criterion]
        Let \(f: [a, b] \to \RR\) be bounded. Then \(f\) is Riemann integrable if and only if, for any \(\eps > 0\), there exists a partition \(\mc P\) of \([a, b]\) such that
        \begin{align*}
            U(f, \mc P) - L(f, \mc P) < \eps
        \end{align*}
    \end{proposition}

    \begin{proof}
        \begin{itemize}
            \item[\(\impliedby\)] We have, for any \(\eps > 0\),
            \begin{align*}
                0 \leq I^*(f) - I_*(f) \leq U(f, \mc P) - L(f, \mc P) < \eps
            \end{align*}
            \item[\(\implies\)] From the definition of \(\sup\) and \(\inf\), we can find partitions \(\mc P\) and \(\mc P'\) such that
            \begin{align*}
                U(f, \mc P) &\leq I^*(f) + \frac{\eps} 2 \\
                L(f, \mc P) &\geq I_*(f) - \frac{\eps} 2
            \end{align*}
            which implies \(U(f, \mc P) - L(f, \mc P') \leq \eps\). To conclude, we need to show this for a single partition. Let \(\mc P'' = \mc P \cup \mc P'\), then
            \begin{align*}
                U(f, \mc P'') - L(f, \mc P'') \leq U(f, \mc P) - L(f, \mc P') \leq \eps < 2\eps
            \end{align*}
        \end{itemize}
    \end{proof}
    
    \begin{remark*}
        As an exericse, one can show that this is equivalent to a sequential characterisation -- finding a sequence of partitions \((\mc P_n)\) such that
        \begin{align*}
            U(f, \mc P_n) - L(f, \mc P_n) \to 0
        \end{align*}
        as \(n \to \infty\).        
    \end{remark*}
    
    \begin{example*}
        Let us attempt to integrate \(f(x) = x\) on \([0, 1]\). Let \(\mc P_n = \{0, \frac 1 n, \dots, \frac {n-1} n, 1\}\). Then
        \begin{align*}
            U(f, \mc P_n) &= \sum_{j=1}^n \frac 1 n \sup_{\left[\frac{j-1}n, \frac jn\right]} x \\
            &= \frac 1 {n^2} \sum_{j=1}^n \sup_{[j-1, j]} x = \frac 1 {n^2} \sum_{j=1}^n j \\
            &= \frac 1 2 \left(1 + \frac 1 n\right)
        \end{align*}
        and 
        \begin{align*}
            L(f, \mc P_n) &= \sum_{j=1}^n \frac 1 n \inf_{\left[\frac{j-1}n, \frac jn\right]} x \\
            &= \frac 1 {n^2} \sum_{j=1}^n \inf_{[j-1, j]} x = \frac 1 {n^2} \sum_{j=1}^n j-1 \\
            &= \frac 1 2 \left(1 - \frac 1 n\right)
        \end{align*}
        Hence
        \begin{align*}
            I_*(f) = \sup_{\mc P}L(f, \mc P) \geq \sup_n L(f, \mc P_n) &= \frac 1 2 \\
            I^*(f) = \inf_{\mc P}U(f, \mc P) \leq \inf_n U(f, \mc P_n) &= \frac 1 2
        \end{align*}
        so \(I_*(f) \geq I^*(f)\), so \(f\) is integrable as required, with \(\int_0^1 x = \frac 1 2\).
        
        Alternatively, with the Riemann integration criterion,
        \begin{align*}
            U(f, \mc P_n) - L(f, \mc P_n) \leq \frac 1 n \to 0
        \end{align*}
        as \(n \to \infty\) (but this doesn't give us the value of the integral).
    \end{example*}
    
    \begin{example*}
        Let us consider a discontinuous but integrable function, namely
        \begin{align*}
            f(x) = \begin{cases}
                0 & x \leq \frac 1 2 \\
                1 & x > \frac 1 2
            \end{cases}
        \end{align*}
        on \([0, 1]\). Then take \(\mc P_n\) as before for \(n \geq 3\). So
        \begin{align*}
            \sup_{[x_{n-1},x_n]}f - \inf_{[x_{n-1},x_n]}f = \begin{cases}
                1 & \frac 1 2 \in [x_{n-1}, x_n] \\
                0 & text{o/w}
            \end{cases}
        \end{align*}
        so that
        \begin{align*}
            U(f, \mc P_n) - L(f, \mc P_n) = \frac 1 n \to 0
        \end{align*}
        as \(n \to \infty\).
    \end{example*}
    
    We would like now to generalise the above examples to get whole classes of functions that are integrable.

    \begin{proposition}[Monotone \(\implies\) integrable]
        Let \(f: [a, b] \to \RR\) be bounded and monotone. Then \(f\) is integrable.
    \end{proposition}

    \lecture{Lecture 16 -- 28/02/26}
    
    \begin{proof}
        Since \(f\) is monotone, \(f\) is bounded by \(\max\left(|f(a)|, |f(b)|\right)\). Without loss of generality, suppose \(f\) is increasing, so
        \begin{align*}
            U(f, \mc P) - L(f, \mc P) &= \sum_{j=1}^n \left(\sup_{I_j} f - \inf_{I_j} f\right)|I_j| \\
            &= \sum_{j=1}^n \left(f(x_j)-f(x_{j-1})\right)|I_j|
        \end{align*}
        for any partition \(\mc P = \{x_0, \dots, x_n\}\). Now, set
        \begin{align*}
            \mc P_n = \left\{a, a + \frac 1 n (b-a), \dots, a + \frac{n-1}{n}(b-a), b\right\}
        \end{align*}
        so that \(|I_j| = \frac{b-a}{n}\). Then
        \begin{align*}
            U(f, \mc P_n) - L(f, \mc P_n) &= \frac{b-a}n \sum_{j=1}^n \frac{b-a}{n} \left[f\left(a + \frac j n (b-a)\right) - f\left(a + \frac{j-1} n (b-a)\right)\right] \\
            &= \frac{b-a}{n}(f(b)-f(a))
        \end{align*}
        since the sum telescopes. But this clearly goes to 0 as \(n \to \infty\).
    \end{proof}

    \begin{proposition}[Continuous \(\implies\) integrable]
        Let \(f: [a, b] \to \RR\) be continuous. Then \(f\) is integrable.
    \end{proposition}
    
    \begin{proof}
        Since \([a, b]\) is a closed bounded set, \(f\) is bounded by \cref{evt} (EVT).
        
        We will show the contrapositive. Suppose \(f\) is not integrable, so there exists \(\eps > 0\) such that, for any \(\mc P\),
        \begin{align*}
            U(f, \mc P) - L(f, \mc P) \geq \eps
        \end{align*}
        Now observe
        \begin{align*}
            \eps \leq U(f, \mc P) - L(f, \mc P) = \sum_{j=1}^n |I_j| \left(\sup_{I_j} f - \inf_{I_j} f\right) \leq \sum_{k=1}^n |I_j| \max_{1 \leq j \leq n} \left(\sup_{I_j} f - \inf_{I_j} f\right)
        \end{align*}
        Now the max is independent of the sum, so we obtain
        \begin{align*}
            \max_{1 \leq j \leq n} \left(\sup_{I_j} f - \inf_{I_j} f\right) \geq \frac{\eps}{b-a}
        \end{align*}
        so, in particular, we obtain a particular \(j\) such that
        \begin{align*}
            \sup_{I_j} f - \inf_{I_j} f \geq \frac{\eps}{b-a}
        \end{align*}
        Since \(I_j\) is closed and bounded, we apply \cref{evt} (EVT) to obtain \(y\), \(z \in [x_{j-1}, x_j]\) such that
        \begin{align*}
            f(y) - f(z) \geq \frac{\eps}{b-a}
        \end{align*}
        Applying this result to \(\mc P_n = \left\{a, a + \frac 1 n (b-a), \dots, a + \frac{n-1}{n}(b-a), b\right\}\) as used previously, we obtain for each \(n\) values of \(y_n\) and \(z_n\) satisfying
        \begin{align*}
            |y_n - z_n| < \frac{b-a}{n} \qquad\qquad |f(y_n)-f(z_n)| \geq \frac{\eps}{b-a}
        \end{align*}
        Now apply \cref{bw} to \((y_n)\) to obtain \((y_{n_k})\) convergent to \(y \in [a, b]\) -- then,
        \begin{align*}
            |z_{n_k}-y_{n_k}| < \frac{b-a}{n_k} \to 0
        \end{align*}
        as \(k \to infty\), so \(z_{n_k} \to y\) also. However, since \(|f(z_{n_k})-f(y_{n_k})| \geq \frac{\eps}{b-a}\), we conclude \(f(z_{n_k})\) and \(f(y_{n_k})\) do not have the same limit, so at least one must not tend to \(f(y)\). Hence \(f\) is not sequentially continuous, and so it is not continuous either.       
    \end{proof}

    \begin{lemma}
        \label{finite-diff-integrable}
        Let \(f\), \(g: [a, b] \to \RR\) be bounded. Suppose \(g\) is integrable, and 
        \begin{align*}
            \{x \in [a, b]: f(x) \neq g(x)\}
        \end{align*}
        is finite. Then \(f\) is integrable and \(\int_a^b f = \int_a^b g\).
    \end{lemma}
    
    \begin{proof}
        Let \(M = \sup_{[a, b]} |f|\). Fix \(\eps > 0\), then there exists a partition \(\mc P\) of \([a, b]\) such that \(U(g, \mc P) - L(g, \mc P) < \eps\). Let \(\{x \in [a, b]: f(x) \neq g(x)\} = {z_1, \dots, z_N}\).

        For each \(k\), \(1 \leq k \leq N\), let \(J_k = [z_k-r_k, z_k+r_k]\), where the \(r_k\) satisfy
        \begin{align*}
            2\sum_{n=1}^N r_k < \eps
        \end{align*}
        so that \(\sum_{k=1}^N |J_k| < \eps\). Let \(\mc P' = \mc P \cup \{z_1 \pm r_1, \dots, z_N \pm r_N\}\), and let \(J = \bigcup J_k\). Clearly, \(U(g, \mc P') - L(g, \mc P') \leq U(g, \mc P) - L(g, \mc P) < \eps\). Now,
        \begin{align*}
            U(f, \mc P') - L(f, \mc P') &= \sum_{j} |I_j| \left(\sup_{I_j} f - \inf_{I_j} f\right) \\
            &= \sum_{I_j \subset J} |I_j| \left(\sup_{I_j} f - \inf_{I_j} f\right) + \sum_{I_j \not\subset J} |I_j| \left(\sup_{I_j} f - \inf_{I_j} f\right)
        \end{align*}
        For the second sum, we have
        \begin{align*}
            \sum_{I_j \not\subset J} |I_j| \left(\sup_{I_j} f - \inf_{I_j} f\right) &= \sum_{I_j \not\subset J} |I_j| \left(\sup_{I_j} g - \inf_{I_j} g\right) \\
            &\leq \sum_j |I_j| \left(\sup_{I_j} g - \inf_{I_j} g\right) \\
            &= U(g, \mc P') - L(g, \mc P') < \eps
        \end{align*}
        For the first sum, we instead know \(\sup_{I_j} f - \inf_{I_j} f \leq 2M\), so that
        \begin{align*}
            \sum_{I_j \subset J} \left(\sup_{I_j} f - \inf_{I_j} f\right) &\leq 2M \sum_{I_j \subset J} |I_j| \\
            &\leq 2M |J| \\
            &< 2M\eps
        \end{align*}
        Thus \(U(f, \mc P') - L(f, \mc P') < (2M+1)\eps\), so \(f\) is integrable. It remains to show \(\int_a^b f = \int_a^b g\). Observe
        \begin{align*}
            L(f, \mc P') \leq \int_a^b f \leq U(f, \mc P') \qquad \qquad L(g, \mc P') \leq \int_a^b g \leq U(g, \mc P')
        \end{align*}
        Hence, again taking \(\mc P\) such that \(U(g, \mc P) - L(g, \mc P) < \eps\),
        \begin{align*}
            \int_a^b f - \int_a^b g &\leq U(f, \mc P') - U(g, \mc P') + \underbrace{U(g, \mc P') - L(g, \mc P')}_{< \eps}
        \end{align*}
        Looking at the first term, we can do basically the exact same thing as above -- but, for completeness,
        \begin{align*}
            U(f, \mc P') - U(g, \mc P') &= \sum_j |I_j| \left(\sup_{I_j}f - \sup_{I_j}g\right) \\
            &= \sum_{I_j \subset J} |I_j| \left(\sup_{I_j}f - \sup_{I_j}g\right) \\
            &\leq \sum_{I_j \subset J} |I_j| \sup_{I_j}f 
            &\leq M|J| < M\eps
        \end{align*}
        A similar argument could be employed to bound the difference below, so that
        \begin{align*}
            \left|\int_a^b f - \int_a^b g\right| < (M+1)\eps 
        \end{align*}
        for any \(\eps > 0\), so, taking the inf over \(\eps\), the LHS is 0 as required.
    \end{proof}
    
    \lecture{Lecture 17 -- 03/03/26}

    \begin{proposition}[Additivity of integrals]
        \label{additivity-integrals}
        Let \(f: [a, b] \to \RR\) and \(c \in (a, b)\). Then \(f\) is integrable if and only if \(\left.f\right|_{[a,c]}\) and \(\left.f\right|_{[c,b]}\) are integrable, and furthermore
        \begin{align*}
            \int_a^b f = \int_a^c f + \int_c^b f
        \end{align*}
    \end{proposition}
    
    \begin{proof}
        Let \(I_L = [a, c]\) and \(I_R = [c, b]\), and \(f_L = \left.f\right|_{I_L}\), \(f_R = \left.f\right|_{I_R}\).
        \begin{itemize}
            \item[\(\implies\)] Fix \(\eps > 0\). Then, there exists a partition \(\mc P\) of \([a, b]\) such that \(U(f, \mc P) - L(f, \mc P) < \eps\). Without loss of generality, \(c = x_l\) for some \(l \in \{0, \dots, n\}\) (else add \(c\) to \(\mc P\) -- this can only make \(U(f, \mc P) - L(f, \mc P)\) smaller). Then \(\mc P = \mc P_L \cup \mc P_R\) where \(\mc P_L = \{a = x_0, \dots, x_l = c\}\) and \(\mc P_R = \{c = x_l, \dots, x_n = b\}\) are partitions of \([a, c]\) and \([c, b]\) respectively. Furthermore,
            \begin{align*}
                U(f, \mc P) &= U(f_L, \mc P_L) + U(f_R, \mc P_R) \\
                L(f, \mc P) &= L(f_L, \mc P_L) + L(f_R, \mc P_R)
            \end{align*}
            so that
            \begin{align*}
                \left(U(f_L, \mc P_L)-L(f_L, \mc P_L)\right) + \left(U(f_R, \mc P_R)-L(f_R, \mc P_R)\right) < \eps
            \end{align*}
            Since both brackets are \(\geq 0\), they must each be \(< \eps\). Hence \(f_L\) and \(f_R\) are both integrable.

            \item[\(\impliedby\)] Fix \(\eps > 0\), so there exist \(\mc P_L\), \(\mc P_R\) partitions of \(I_L\), \(I_R\) respectively such that
            \begin{align*}
                U(f_R, \mc P_R) - L(f_R, \mc P_R) &< \eps
                U(f_L, \mc P_L) - L(f_L, \mc P_L) &< \eps
            \end{align*}
            so \(U(f, \mc P) - L(f, \mc P) < 2\eps\), since \(\mc P = \mc P_L \cup \mc P_R\) is a partition of \([a, b]\). 
        \end{itemize}
        Furthermore,
        \begin{align*}
            \int_a^b f &\geq L(f_L, \mc P_L) + L(f_R, \mc P_R) \\
            &\geq \int_a^c f_L + \int_c^b f_R - 2\eps
        \end{align*}
        and, similarly,
        \begin{align*}
            \int_a^b f \leq \int_a^c f_L + \int_c^b f_R + 2\eps
        \end{align*}
        so
        \begin{align*}
            \left|\int_a^b f - \left(\int_a^c f_L + \int_c^b f_R\right)\right| < 2\eps
        \end{align*}
        Since \(\eps\) is arbitrary, this is what we wanted.        
    \end{proof}

    \begin{definition}[Piecewise continuity]
        Let \(f: [a, b] \to \RR\). Then \(f\) is \emph{piecewise continuous} if there exists a partition \(\mc P\) such that \(\left.f\right|_{(x_{j-1},x_j)}\) is continuous for each \(j\), and has a finite limit as \(x \to x_{j-1}\) from above and \(x \to x_j\) from below.
    \end{definition}
    
    \begin{proposition}[Piecewise continuous \(\implies\) integrable]
        Let \(f: [a, b] \to \RR\) be piecewise continuous with \(\mc P\) the relevant partition. Then \(f\) is integrable, and
        \begin{align*}
            \int_a^b f(x) \D x = \sum_{j=1}^n \int_{x_{j-1}}^{x_j} \bar{f}_j(x) \D x
        \end{align*}
        where
        \begin{align*}
            \bar f_j(x) = \begin{cases}
                f(x) & x \in (x_{j-1}, x_j) \\
                \lim_{x \to x_j^-}f(x) & x = x_j \\
                \lim_{x \to x_{j-1}^+}f(x) & x = x_{j-1}
            \end{cases}
        \end{align*}
    \end{proposition}

    \begin{proof}
        \Cref{finite-diff-integrable} gives that \(\left.f\right|_{[x_{j-1},x_j]}\) is integrable and \(\int_{x_{j-1}}^{x_j} f = \int_{x_{j-1}}^{x_j} \bar f_j\). The result follows from \cref{additivity-integrals}.
    \end{proof}

    Continuous, piecewise continuous and monotone functions are all integrable, as are functions made from adding these. Are these all the Riemann integrable functions? No.

    \begin{example*}[Thomae/popcorn function]
        Define
        \begin{align*}
            f: [0, 1] &\to \RR \\
            x &\mapsto \begin{cases}
                \frac 1 q & x = \frac p q,\, p,\, q \text{ coprime positive integers} \\
                0 & x \in \RR \setminus \QQ
            \end{cases}
        \end{align*}
        Since \(\RR \setminus \QQ\) is dense in \(\RR\), \(L(f, \mc P) = 0\) for any partition \(\mc P\) of \([0, 1]\), so \(I_*(f) = 0\). We claim \(f\) is integrable, so, for any \(\eps > 0\), it is necessary to find a partition \(\mc P\) of \([0, 1]\) such that \(U(f, \mc P) < \eps\).

        Pick \(N \in \NN\) such that \(N > \frac 1 \eps\). Set \(X_N = \{x \in [0, 1]: f(x) \geq \frac 1 N\} \subseteq \{\frac p q: 1 \leq q \leq N, 0 \leq p \leq q\} = \{y_1, \dots, y_M\}\), which is clearly finite. Define \(\mc P\) such that
        \begin{itemize}
            \item each \(y_k\) is in some subinterval of \(\mc P\);
            \item each subinterval has length \(< \frac{\eps} M\).
        \end{itemize}
        Now
        \begin{align*}
            U(f, \mc P) &= \sum_{\substack{I \in \mc P \\ I \cap X_N \neq \emptyset}} |I| \underbrace{\sup_I f}_{\leq 1} + \sum_{\substack{I \in \mc P \\ I \cap X_N = \emptyset}} |I| \underbrace{\sup_I f}_{\leq \frac 1 N} \\ 
            &\leq \sum_{\substack{I \in \mc P \\ I \cap X_N \neq \emptyset}} |I| + \frac 1 N \sum_{I \in \mc P} |I| \\
            &\leq M \cdot \frac{\eps} M + \eps < 2\eps
        \end{align*}
    \end{example*}
    
    \begin{proposition}[Countable discontinuities \(\implies\) integrable]
        Let \(f\) bounded \(: [a, b] \to \RR\), and let \(D\) be the set of discontinuities of \(f\). If \(D\) is finite or countable, \(f\) is Riemann-integrable.
    \end{proposition}
    
    \begin{remark*}
        The proofs of these results are on sheet 3 and are too hard for the course, respectively.
    \end{remark*}
    
    \begin{remark*}
        Even now there are still more integrable functions we have not categorised. For example, the indicator function of the Cantor set is integrable but has uncountably many discontinuities.
    \end{remark*}
    
    Let us have some boring lemmas.

    \begin{lemma}
        Let \(I\) be a closed bounded interval, and let \(f, g\) bounded \(: I \to \RR\). Then
        \begin{enumerate}
            \item \(f(x) \leq g(x)\) for all \(x \in I \implies\) sup and inf preserve the inequality;
            \item \(\sup_I (-f) = -\inf_I f\);
            \item for \(\lambda > 0\), \(\sup_I(\lambda f) = \lambda \sup_I f\);
            \item \(\sup_I(f+g) \leq \sup_I f + \sup_I g\), and \(\inf_I(f+g) \geq \inf_I f + \inf_I g\), without equality necessarily holding;
            \item \(\sup_I |f| - \inf_I |f| \leq \sup_I f - \inf_I f\);
            \item \(\sup_I f^2 - \inf_I f^2 \leq 2 \sup_I |f|\left(\sup_I f - \inf_I f\right)\).
        \end{enumerate}
    \end{lemma}
    
    \lecture{Lecture 18 -- 05/03/26}

    \begin{proof}
        We leave 1 to 4 and part of 5 as an exercise. Then
        \begin{enumerate}
            \item[5.] If \(f \geq 0\) or \(f \leq 0\) on all of \(I\), the reuslt is immediate, or follows from 2. Otherwise, if \(\inf_I f < 0 < \sup_I f\), then
            \begin{align*}
                \text{LHS} &\leq \sup_I |f| \leq \max(\sup_I f, \sup_I (-f)) \\
                &\leq \sup_I f + \sup_I(-f) = \sup_I f - \inf_I f
            \end{align*}
            \item [6.] We have
            \begin{align*}
                f(x)^2 - f(y)^2 &= (f(x)-f(y))(f(x)+f(y)) \\
                \sup_I f^2 - \inf_I f^2 &\leq \sup_{x \in I} \inf_{y \in I} \left((f(x)-f(y))(f(x)+f(y))\right) \\
                &\leq 2 \sup_I |f| \left(\sup_I f - \inf_I f\right)
            \end{align*}
        \end{enumerate}
    \end{proof}

    \begin{lemma}
        Let \(f\), \(g: [a, b] \to \RR\) be integrable. Then
        \begin{enumerate}
            \item \(f(x) \leq g(x)\) for all \(x \implies \int_a^b f \leq \int_a^b g\);
            \item[2, 3.] for \(\lambda \in \RR\), \(\int_a^b \lambda f = \lambda \int_a^b f\);
            \item[4] \(f+g\) is integrable, and \(\int_a^b (f+g) = \int_a^b f + \int_a^b g\);
            \item[5] \(|f|\) is integrable, and \(\int_a^b |f| \geq \left|\int_a^b f\right|\);
            \item[6] \(fg\) is integrable.
        \end{enumerate}
    \end{lemma}
    
    \begin{proof}
        Again, we leave 1 to 4 as exercises.
        \begin{enumerate}
            \item[5.] For any partition \(\mc P\) of \([a, b]\),
            \begin{align*}
                U(|f|, \mc P) - L(|f|, \mc P) \leq U(f, \mc P) - L(f, \mc P)
            \end{align*}
            so, if \(f\) is integrable, then \(|f|\) is integrable. Now
            \begin{align*}
                -|f(x)| \leq f(x) \leq |f(x)|
            \end{align*}
            so, by 1,
            \begin{align*}
                -\int_a^b |f(x)| \D x \leq \int_a^b f(x) \D x \leq \int_a^b |f(x)| \D x \\
                \implies \int_a^b |f(x)| \D x \geq \left|\int_a^b f(x) \D x\right|
            \end{align*}
            as required.
            \item[6.] Observe
            \begin{align*}
                f(x)g(x) = \frac 1 4 \left((f(x)+g(x))^2-(f(x)-g(x))^2\right)
            \end{align*}
            Hence, it is only necessary to show that, if \(f\) is integrable, \(f^2\) is integrable. For any partition \(\mc P\) of \([a, b]\), we have
            \begin{align*}
                U(f^2, \mc P) - L(f^2, \mc P) \leq 2 \sup_I |f| \left(U(f, \mc P) - L(f, \mc P)\right)
            \end{align*}
            The result follows.
        \end{enumerate}
    \end{proof}
    
    \subsection{Integration and differentiation}

    We all know that integration and differentiation are famously opposites -- but why?

    \begin{definition}[Indefinite integral]
        The \emph{indefinite integral} of a function \(f: [a, b] \to \RR\) is
        \begin{align*}
            F: [a, b] &\to \RR \\
            x &\mapsto \int_a^x f(t) \D t
        \end{align*}
    \end{definition}
    
    \begin{proposition}[Integrals are continuous]
        Let \(f: [a, b] \to \RR\) be Riemann integrable, and let \(F\) be its indefinite integral. Then \(F\) is continuous.
    \end{proposition}
    
    \begin{proof}
        Observe
        \begin{align*}
            |F(x+h)-F(x)| &= \left|\int_a^{x+h} f(t) \D t - \int_a^x f(t) \D t \right| \\
            &= \left|\int_a^{x+h} f(t) \D t\right| \\
            &\leq \int_x^{x+h} |f(t)| \D t \\
            &\leq \sup_{[a,b]} f \int_x^{x+h} 1 \D t
        \end{align*}
        which goes to 0 as \(h \to 0\).
    \end{proof}
    
    \begin{theorem}[Fundamental theorem of calculus, part 1]
        \label{ftc-1}
        If \(f: [a, b] \to \RR\) is Riemann integrable and continuous at \(x_0\), then its indefinite integral \(F\) is differentiable at \(x_0\). Furthermore,
        \begin{align*}
            F'(x_0) = \left.\frac{d}{dx} \int_a^x f(t) \D t\right|_{x_0} = f(x_0)
        \end{align*}
    \end{theorem}
    
    \begin{proof}
        Let
        \begin{align*}
            \xi(h) = \frac{F(x_0+h)-F(x_0)-hf(x_0)}{|h|}
        \end{align*}
        and we will show \(\xi(h) \to 0\) as \(h \to 0\). Observe
        \begin{align*}
            |F(x_0+h)-F(x_0)-hf(x_0)| &= \left|\int_x^{x_0+h} f(t) \D t - hf(x_0)\right| \\
            &= \left|\int_x^{x_0+h} (f(t)-f(x_0)) \D t\right| \\
            &\leq \int_x^{x_0+h} |f(t)-f(x_0)| \D t \\
            &= \sup_{t \in [0, h]}|f(x_0+t)-f(x_0)| \cdot \underbrace{\int_x^{x_0+h} \D t}_h
        \end{align*}
        Hence
        \begin{align*}
            \xi(h) \leq \sup_{t \in [0, h]}|f(x_0+t)-f(x_0)|
        \end{align*}
        which goes to 0 by continuity as required. 
    \end{proof}
    
    \begin{example*}
        Take \([a, b] = [-1, 1]\), and
        \begin{align*}
            f(x) = \begin{cases}
                -1 & x \leq 0 \\
                1 & x > 0
            \end{cases}
        \end{align*}
        which is piecewise continuous hence integrable, with
        \begin{align*}
            F(x) = \begin{cases}
                -1-x & x \leq 0 \\
                x-1 & x > 0
            \end{cases} = |x| - 1
        \end{align*}
        which is not differentiable at \(x=0\). Hence, continuity is necessary in the theorem.       
    \end{example*}
    
    \begin{corollary}
        If \(f = g'\) continuous on \([a, b]\), then
        \begin{align*}
            F(x) = \int_a^x f(t) \D t = g(x) - g(a)
        \end{align*}
    \end{corollary}
    
    \begin{proof}
        Observe \((F-g)' = 0\) on \([a, b]\), so \(F-g\) is constant on \([a, b]\) by the mean value theorem. Then \(F(x) - g(x) = F(a) - g(a)\), but \(F(a) = 0\), yielding the result.
    \end{proof}
    
    \begin{theorem}[Fundamental theorem of calculus, part 2]
        \label{ftc-2}
        If \(f: [a, b] \to \RR\) is Riemann integrable and there exists a differentiable \(F: [a, b] \to \RR\) such that \(F'(x) = f(x)\), then
        \begin{align*}
            \int_a^b f(t) \D t = \int_a^b F'(t) \D t = F(b) - F(a)
        \end{align*}
    \end{theorem}
    
    \begin{proof}
        By assumption, there exists for any \(\eps > 0\) a partition \(\mc P\) of \([a, b]\) such that \(U(f, \mc P) - L(f, \mc P) < \eps\). Applying the mean value theorem to \(F\) on intervals of this partition, we have
        \begin{align*}
            F(x_j) - F(x_{j-1}) = f(t_j)(x_j-x_{j-1})
        \end{align*}
        for some \(t_j \in (x_{j-1}, x_j)\). Now
        \begin{align*}
            F(b) - F(a) = \sum_{j=1}^n f(t_j)(x_j-x_{j-1}) \in [L(f, \mc P), U(f, \mc P)]
        \end{align*}
        so, taking inf and sup,
        \begin{align*}
            I_*(f) \leq F(b)-F(a) \leq I^*(f)
        \end{align*}
        so \(F(b)-F(a) = \int_a^b f\) as required.
    \end{proof}

    \lecture{Lecture 19 -- 07/03/26}
    
    Let us now derive some important consequences of this result.

    \begin{proposition}[Integration by parts]
        Suppose \(f\), \(g \in C^1([a, b])\). Then
        \begin{align*}
            \int_a^b f'g = f(b)g(b) - f(a)g(a) - \int_a^b fg'
        \end{align*}
    \end{proposition}
    
    \begin{proof}
        The product rule gives \(f'g = (fg)' - fg'\). Integrating this over \([a, b]\) and applying \cref{ftc-2} to evaluate the \((fg')\) integral gives the result.
    \end{proof}
    
    \begin{proposition}[Integration by substitution]
        Let \(f: [a, b] \to \RR\) be continuous, and let \(g: [\alpha, \beta] \to [a, b]\) satisfy \(g \in C^1([\alpha, \beta])\) with \(g(\alpha) = a\), \(g(\beta) = b\). Then
        \begin{align*}
            \int_a^b f(x) \D x = \int_\alpha^\beta f(g(t)) g'(t) \D t
        \end{align*}
    \end{proposition}
    
    \begin{proof}
        Let \(F(x) = \int_a^x f(t) \D t\), then \(F: [a, b] \to \RR\) is well-defined and differentiable by \cref{ftc-1}. Set \(h = F \circ g: [\alpha, \beta] \to \RR\). Then \(h\) is differentiable, with
        \begin{align*}
            h'(t) = F'(g(t))g(t) = f(g(t))g'(t)
        \end{align*}
        by the chain rule and \cref{ftc-1}. Hence, applying \(\cref{ftc-2}\),
        \begin{align*}
            \int_a^b f(x) \D x &= F(b) - F(a) = h(\beta) - h(\alpha) \\
            &= \int_\alpha^\beta h'(t) \D t \\
            &= \int_\alpha^\beta f(g(t)) g'(t) \D t
        \end{align*}
        by \cref{ftc-2}
    \end{proof}
    
    \begin{theorem}[Taylor's theorem, integral remainder]
        Let \(f \in C^n([a, a+h])\), and let \(R_{n,f,a}\) be as before. Then
        \begin{align*}
            R_{n,f,a}(h) &= \frac{h^n}{(n-1)!} \int_0^1 (1-t)^{n-1}f^{(n)}(a+th) \D t \\
            &= \frac 1 {(n-1)!} \int_0^h (h-u)^{n-1} f^{(n)}(a+u) \D u
        \end{align*}
    \end{theorem}
    
    \begin{remark*}
        Letting \(M_n = \sup_{[a,a+h]}|f^(n)|\) as before, we know
        \begin{align*}
            \left|R_{n,f,a}(h)\right| &\leq \frac{h^n}{n!} \int_0^1 \left|(1-t)^{h-1} f^{(n)}(a+th)\right| \D t \\
            &\leq \frac{h^n}{n!} M_n \int_0^1 1 \D t
        \end{align*}
        so \(R_{n,f,a}(h) \to 0\) as \(h \to 0\) for \(n\) fixed. Furthermore, if we know \(\sup_{n \geq 0} M_n = M < \infty\), we would obtain
        \begin{align*}
            |R_{n,f,a}(h)| \leq \frac{Mh^n}{h!} \to 0
        \end{align*}
        as \(n \to \infty\) for all \(|h| < 1\), so \(f\) analytic around \(a\).       
    \end{remark*}
    
    \begin{proof}
        We will use integration by parts. Observe
        \begin{align*}
            &\; \frac{1}{(n-1)!} \int_0^h (h-u)^{n-1} f^{(n)}(a+u) \D u \\
            =&\; h^{n-1} \frac{f^{(n-1)}(a)}{(n-1)!} + \frac{1}{(n-2)!} \int_0^h (h-u)^{n-2} f^{(n-1)}(a+u) \D u \\
            =&; \dots = -\sum_{k=1}^{n-1} \frac 1 {n!} f^{(k)}(a) h^k + \underbrace{\int_0^h f'(a+u) \D u}_{f(a+h)-f(a)} \\
            =&; R_{n,f,a}(h)
        \end{align*}
        as required.        
    \end{proof}
    
    \begin{proposition}[Cauchy mean value theorem for integrals]
        Let \(f\), \(g: [a, b] \to \RR\) continuous, \(g(x) \neq 0\) for any \(x \in [a, b]\). Then there exists \(c \in (a, b)\) such that
        \begin{align*}
            \int_a^b f(x)g(x) \D x = f(c) \int_a^b g(x) \D x
        \end{align*}
    \end{proposition}
    
    \begin{proof}
        Let us apply Cauchy MVT to \(F(x) = \int_a^x f(t)g(t) \D t\), \(G(x) = \int_a^x g(t) \D t\), which are differentiable by \cref{ftc-1}. We obtain \(c \in (a, b)\) such that
        \begin{align*}
            g(c) \int_a^b f(t)g(t) \D t = (F(b)-F(a))G'(c) = F'(c)(G(b)-G(a)) = f(c)g(c) \int_a^b g(t) \D t
        \end{align*}
        so the result follows after division by \(g(c)\).        
    \end{proof}
    
    We can now give alternative proofs of the Lagrange (\cref{taylor-lagrange}) and Cauchy (\cref{taylor-cauchy}) remainders for Taylor's theorem.

    \begin{proof}[Lagrange remainder (\emph{assuming \(f^{(n)}\) continuous})]
        Let \(g(t) = (1-t)^{n-1}\), then there exists \(\theta \in (0, 1)\) such that
        \begin{align*}
            R_{n,f,a} &= \frac{h^n}{(h-1)!} f^{(n)}(a + \theta h) \int_0^1 (1-t)^{n-1} \D t \\
            &= \frac{h^n}{n!} f^{(n)}(a+\theta h)
        \end{align*}
    \end{proof}
    
    \begin{proof}[Cauchy remainder (\emph{asumming \(f^{(n)}\) continuous})]
        Take \(g=1\). Then
        \begin{align*}
            R_{n,f,a}(h) &= \frac{h^n}{(n-1)!} (1-\theta)^{n-1} f^{(n)}(a+\theta h) \int_0^1 \D t \\
            &= \frac{h^n}{(n-1)!} (1-\theta)^{n-1} f^{(n)}(a+\theta h)
        \end{align*}
        as required.
    \end{proof}
    
    \subsection{Improper integrals}

    We know how to integrate bounded functions on bounded domains, but what about functions with unbounded image or domain? Let's start with unbounded domains.

    \begin{definition}[Improper integrals on unbounded domain]
        Suppose \(f: [a, \infty) \to \RR\) is integrable on \([a, R]\) for every \(a \leq R < \infty\), and set
        \begin{align*}
            F: [a, \infty) &\to \RR \\
            R &\mapsto \int_a^R f(x) \D x
        \end{align*}
        Then we say \(\int_a^\infty f(x) \D x\) exists/converges if \(\lim_{R \to \infty} = L \in \RR\) exists and is finite, in which case we define
        \begin{align*}
            \int_a^\infty f(x) \D x = L
        \end{align*}
        Otherwise, we say \(\int_a^\infty f(x) \D x\) does not exist/does not converge.
        
        If \(f: \RR \to \RR\) is such that \(\int_a^\infty f(x) \D x = L_1\) and \(\int_{-\infty}^a f(x) \D x = L_2\), then we say \(\int_{-\infty}^\infty f(x) \D x\) exists, and we set
        \begin{align*}
            \int_{-\infty}^\infty f(x) \D x = L_1 + L_2 = \lim_{R \to \infty} \int_a^R f(x) \D x + \lim_{r \to \infty} \int_{-r}^a f(x) \D x
        \end{align*}
    \end{definition}
    
    \begin{remark*}
        This is not the same thing as
        \begin{align*}
            \lim_{R \to \infty} \int_{-R}^R f(x) \D x
        \end{align*}
        which is called the \emph{principal value} of the integral. Indeed, for the integral to exist, the limit must exist irrespectively of how quickly we approach \(\infty\) and \(-\infty\) on each side.        
    \end{remark*}
    
    \begin{example*}
        \begin{itemize}
            \item Take
            \begin{align*}
                \int_1^\infty \frac{1}{x^p} \D x = \lim_{R \to \infty} \int_1^R \frac{1}{x^p} \D x
            \end{align*}
            which exists if and only if \(p > 1\).
            \item Take
            \begin{align*}
                \int_2^\infty \frac{1}{x \log^2 x }\D x = \lim_{R \to \infty} \int_2^R \frac{1}{x \log^2 x}\D x = \frac 1{\log 2}
            \end{align*}
        \end{itemize}
    \end{example*}
    
    \begin{proposition}[Comparison test for integrals]
        Suppose \(f\), \(g: [a, \infty)\) are integrable on \([a, R]\) for every \(R < \infty\), and satisfy \(0 \leq f(x) \leq g(x)\) for any \(x \in [a, \infty)\). Then
        \begin{itemize}
            \item \(\int_a^\infty g(x) \D x\) converges \(\implies \int_a^\infty f(x) \D x\) converges, and \(\int_a^\infty f(x) \D x \leq \int_a^\infty g(x) \D x\);
            \item \(\int_a^\infty f(x) \D x\) diverges to \(\infty\) \(\implies \int_a^\infty g(x) \D x\) diverges to \(\infty\).
        \end{itemize}
    \end{proposition}
    
    \begin{proof}
        \begin{itemize}
            \item Let \(F(R) = \int_a^R f(x) \D x\), which is increasing since \(f \geq 0\). It is also bounded, since
            \begin{align*}
                0 \leq F(R) \leq \int_a^R g(x) \D x \leq \int_a^\infty g(x) \D x
            \end{align*}
            Since \(F\) is monotone and bounded, \(L = \sup_{R \geq a} F(R) < \infty\) exists. We claim \(L = \int_a^\infty f(x) \D x\). By definition of sup, for any \(\eps > 0\), there exists \(R_0 \in [a, \infty)\) such that
            \begin{align*}
                L-\eps \leq F(R_0) \leq F(R) \leq L
            \end{align*}
            Now, taking \(R \to \infty\), we obtain
            \begin{align*}
                L-\eps \leq \lim_{R \to \infty} F(R) \leq L
            \end{align*}
            for any \(\eps > 0\), so \(\lim_{R \to \infty} F(R) = L\) as required.
            \item Since \(f \geq 0\), \(\lim_{R \to \infty} \int_a^R f(x) \D x = \infty\). Hence, for any \(L > 0\), there exists \(R\) such that \(F(r) > L\) for all \(r \geq R\). But then
            \begin{align*}
                \int_a^R g(x) \D x \geq F(R) > L
            \end{align*}
            so, taking the limit on both sides,
            \begin{align*}
                \lim_{R \to \infty} \int_a^R g(x) \D x = \infty
            \end{align*}
            as required.
        \end{itemize}
    \end{proof}
    
    \begin{example*}
        Let us show that \(\int_{-\infty}^\infty e^{-x^2} \D x\) converges. Suppose \(x \geq 1\), then \(x^2 \geq x \implies e^{-x^2} \leq e^{-x}\). Hence
        \begin{align*}
            \int_1^\infty e^{-x^2} \D x \leq \int_1^\infty e^{-x} \D x = \left[-e^{-x}\right]_1^{\infty} = \frac 1 e
        \end{align*}
        so
        \begin{align*}
            \int_{-\infty}^\infty e^{-x^2} \D x = 2 \int_0^\infty e^{-x^2} \D x = \underbrace{2 \int_1^\infty e^{-x^2} \D x}_{\leq \frac 2 e} + \underbrace{2 \int_0^1 e^{-x^2} \D x}_{\text{finite}}
        \end{align*}
    \end{example*}
    
    \begin{proposition}[Ratio test for integrals]
        Suppose \(f\), \(g: [a, \infty)\) are integrable on \([a, R]\) for every \(R < \infty\), satisfy \(f\), \(g \geq 0\), and suppose further
        \begin{align*}
            \lim_{x \to \infty} \frac{f(x)}{g(x)} = L \in (0, \infty)
        \end{align*}
        Then \(\int_a^\infty f(x) \D x\) converges if and only if \(\int_a^\infty g(x) \D x\) converges.
    \end{proposition}
    
    \begin{proof}
        Left as an exercise by applying the comparison test.
    \end{proof}
    
    \begin{remark*}
        If \(\lim_{x \to \infty} \frac{f(x)}{g(x)} = 0\), then \(f(x) \leq g(x)\) from some point, so by comparison test from that point onwards
        \begin{align*}
            \int_a^\infty g(x) \D x \text{ convergent } \implies \int_a^\infty f(x) \D x \text{ convergent} 
        \end{align*}
        Likewise, if \(\lim_{x \to \infty} \frac{f(x)}{g(x)} = \infty\), then \(f(x) \geq g(x)\) from some point, so by comparison test from that point onwards
        \begin{align*}
            \int_a^\infty g(x) \D x \text{ divergent } \implies \int_a^\infty f(x) \D x \text{ divergent }
        \end{align*}
    \end{remark*}
    
    \begin{example*}
        \begin{itemize}
            \item Consider again \(e^{-x^2}\). We have \(\lim_{x \to \infty} \frac{e^{-x^2}}{e^{-x}} = 0\), and \(\int_0^\infty e^{-x} \D x = 1 < \infty\), so applying the remark gives \(\int_0^\infty e^{-x^2} \D x < \infty\).
            \item Consider \(\int_1^\infty \frac x{x^4+1} \D x\). We have
            \begin{align*}
                \lim_{x \to \infty} \frac{\frac{x}{x^4+1}}{\frac 1 {x^3}} = 1
            \end{align*}
            so, by the ratio test, since \(\int_1^\infty \frac 1 {x^3} \D x < \infty\), we know \(\int_1^\infty \frac{x}{x^4+1} \D x < \infty\).
        \end{itemize}
    \end{example*}
    
    Indeed, on sheet 4 we will see that the root test and Dirichlet test also have natural generalisations to integrals. The integral test can also be nicely thought of in the context of this chapter.

    \lecture{Lecture \(n+1\)}

    \begin{definition}[Improper integrals with an isolated singularity]
        Let \(f: (a, b] \to \RR\) be integrable on \([a + \delta, b]\) for any \(0 < \delta \leq b-a\). Set 
        \begin{align*}
            F: (0, b-a] &\to \RR
            \delta &\mapsto \int_{a+\delta}^b f(x) \D x
        \end{align*}
        Then, we say \(\int_a^b f(x) \D x\) exists/converges if \(\lim_{\delta \to 0} F(\delta)\) exists and is finite, and we say
        \begin{align*}
            \int_a^b f(x) \D x = \lim_{\delta \to 0} F(\delta)
        \end{align*}
        Otherwise, we say \(\int_a^b f(x) \D x\) does not exist/does not converge.

        If \(f: [B, b] \setminus \{a\}\) is such that \(\int_B^a f(x) \D x = l_1\) and \(\int_a^b f(x) \D x = l_2\), then we say \(\int_B^b f(x) \D x\) exists, and we set
        \begin{align*}
            \int_B^b f(x) \D x = l_1 + l_2 = \lim_{\delta_1 \to 0} \int_B^{a-\delta_1} f(x) \D x + \lim_{\delta_2 \to 0} \int_{a+\delta_2}^b f(x) \D x
        \end{align*}
    \end{definition}
    
    \begin{remark*}
        Analogously to above, this is not the same thing as
        \begin{align*}
            \lim_{\delta \to 0} \left(\int_B^{a-\delta} f(x) \D x + \int_{a+\delta}^b f(x) \D x\right)
        \end{align*}
        as exemplified by, for example, \(f(x) = \frac 1 x\).
    \end{remark*}
    
    \begin{example*}
        \(\int_0^1 \frac 1 {x^p} \D x\) converges if and only if \(p < 1\). Observe
        \begin{align*}
            \int_\delta^1 \frac 1 {x^p} \D x = \begin{cases}
                [\log x]_{\delta}^1 = -\log \delta & p = 1 \\
                \left[\frac{x^{1-p}}{1-p}\right]_\delta^1 = \frac{\delta^{1-p}-1}{p-1} & p \neq 1
            \end{cases}
        \end{align*}
        Taking \(\delta \to 0\), we see that the limit does not exist if \(p \geq 1\), but exists otherwise.
    \end{example*}
    
    Indeed, it turns out that these two types of improper integral are in some strong sense equivalent -- with a suitable substitution, we can always interconvert between them.

    \begin{lemma}
        Let \(f: (a, b] \to \RR\) be Riemann integrable on \([a+\delta, b]\) for any \(\delta \in (0, b-a)\). Choose \(\phi: [c, \infty) \to (a, b]\) a \(C^1\) strictly decreasing bijection with \(\phi(c) = b\), \(\lim_{t \to \infty}\phi(t) = a\). Then
        \begin{align*}
            \int_a^b f(x) \D x \text{ exists } \iff \int_c^\infty f(\phi(t))(-\phi'(t)) \D t \text{ exists}
        \end{align*}
    \end{lemma}
    
    \lecture{Lecture \(n\)}

    \section{?}

    \lecture{Lecture \(n+1\)}

    \lecture{Lecture \(n+2\)}

    \subsection{Exponential and logarithm}
    
    Taylor's theorem with Lagrange remainder says that
    \begin{align*}
        R_{N,0}(x) = e^x - \sum_{n=1}^{N-1} \sum_{x^n}{n!} = e^\xi \frac{x^N}{N!}
    \end{align*}
    For fixed \(x \in \RR\), \(|R_{N,0}(x)| \leq |e^\xi| \frac{|x|^N}{N!}\), which goes to 0 as \(N \to \infty\). We conclude that, for \(x \in \RR\),
    \begin{align*}
        e^x = \sum_{n=0}^\infty \frac{x^n}{n!}
    \end{align*}
    Now, extending to \(\CC\), we will make this definitional.

    \begin{definition}[Exponential function]
        The \emph{exponential function} is given by
        \begin{align*}
            \exp: \CC &\to \CC \\
            z &\mapsto \sum_{n=0}^\infty \frac{z^n}{n!}
        \end{align*}
    \end{definition}
    
    \begin{lemma}
        \(\exp\) is smooth with \(\exp' = \exp\), \(\exp(0) = 1\), and \(\exp(a+b) = \exp a \exp b\).
    \end{lemma}
    
    \begin{proof}
        The first two properties are clear. For the third part, let \(f(z) = \exp(a+b-z)\exp(z)\), which is smooth and hence differentiable, with \(f'(z) = 0\) after applying \(\exp = \exp'\). Hence \(f\) is constant, so \(f(a) = f(0)\) -- we have
        \begin{align*}
            \exp(a+b-a)\exp(a) = \exp(a+b)
        \end{align*}
        as required.       
    \end{proof}
    
    \begin{lemma}
        Restrict \(\exp\) to \(\RR \to \RR\). Then
        \begin{itemize}
            \item \(\exp(x) > 0\) for all \(x \in \RR\);
            \item \(\exp\) is strictly increasing;
            \item \(\exp(x) \to \infty\) as \(x \to \infty\), \(\exp(x) \to 0\) as \(x \to -\infty\);
            \item \(\exp: \RR \to (0, \infty)\) is a bijection.
        \end{itemize}
    \end{lemma}
    
    \begin{proof}
        For \(x < 0\), every term in the series is positive, so \(\exp(x) \geq 1 + x\), so that \(\exp(x) > 0\) and \(\lim_{x \to \infty}\exp(x) = \infty\). On the other hand, for \(x < 0\), \(\exp(-x)\exp(x) = 1\), so \(\exp(-x) > 0 \implies \exp(x) > 0\), and furthermore \(\exp(-x) \to \infty \implies \exp(x) \to 0^+\).

        That \(\exp\) is strictly increasing follows from \(\exp' = \exp > 0\). From this it follows that \(\exp\) is injective. Now, given
        \begin{align*}
            y \in (0, \infty) = \left(\lim_{x \to -\infty}\exp(x), \lim_{x \to \infty}\exp(x)\right)
        \end{align*}
        Hence, there are reals \(a\), \(b\) with \(\exp(a) < y < \exp(b)\). Applying the intermediate value theorem gives surjectivity and hence bijectivity as required.               
    \end{proof}
    
    Since \(\exp: \RR \to (0, \infty)\) is bijective, it must have an inverse.
    
    \begin{definition}[Natural logarithm]
        The \emph{natural logarithm} is defined as follows:
        \begin{align*}
            \log: (0, \infty) &\to \RR \\
            x &\mapsto \exp^{-1}(x)
        \end{align*}
    \end{definition}
    
    \begin{lemma}
        We have that
        \begin{itemize}
            \item \(\log\) is smooth and monotone, with \(\log'(y) = \frac 1 y\) and \(\log y = \int_1^y \frac 1 t \D t\);
            \item \(\log(yz) = \log y + \log z\);
            \item \(\log(y) \to \infty\) as \(y \to \infty\), \(\log(y) \to -\infty\) as \(y \to 0^+\).
        \end{itemize}
    \end{lemma}
    
    \begin{proof}
        The inverse function theorem says that \(\log\) is smooth, monotone and has derivative \(\frac 1 y\). Then, FTC gives
        \begin{align*}
            \int_1^y \frac 1 t \D t = \log y - \log 1 = \log y
        \end{align*}
        since \(\log y = \log(\exp(0)) = 0\).        

        For the second part, we have \(\log(yz) = \log(\exp(u)\exp(v)) = \log(\exp(u+v)) = u + v = \log(y) + \log(z)\).

        The last part is left as an exercise.
    \end{proof}
    
    \begin{remark*}
        Observe
        \begin{align*}
            \log(1+y) &= \int_1^y \frac 1{1+t} \D t \\
            &= \int_1^y \sum_{n=0}^\infty (-1)^n t^n \D t
        \end{align*}
        if \(|t| \leq |y| < 1\). Since the inner power series has radius of convergence 1, and we are already restricted to \(|t| < 1\), we can integrate term-by-term to obtain
        \begin{align*}
            \log(1+y) = \sum_{n=0}^\infty \frac{(-1)^n}{n+1}y^{n+1}
        \end{align*}
        By taking limits as \(y \to 1\), we can show
        \begin{align*}
            \sum_{n=0}^\infty \frac{(-1)^{n-1}}{n} = \log 2
        \end{align*}
    \end{remark*}
    
    \begin{definition}[Exponentiation]
        Let \(x > 0\), then we define
        \begin{align*}
            x^\alpha := \exp(\alpha \log x)
        \end{align*}
    \end{definition}
    
    \begin{lemma}
        Let \(x\), \(y > 0\), \(\alpha\), \(\beta \in \RR\). Then
        \begin{itemize}
            \item \((xy)^\alpha = x^\alpha y^\alpha\);
            \item \(x^{\alpha+\beta} = x^\alpha x^\beta\);
            \item \((x^\beta)^\alpha = x^{\alpha\beta}\);
            \item \(x^1 = x\), \(x^0 = 1\);
        \end{itemize}
    \end{lemma}
    
    \begin{proof}
        These results mostly follow easily from the definitions. We will check the third -- we have
        \begin{align*}
            (x^\beta)^\alpha &= \exp(\alpha \log(\exp(\beta \log(x)))) \\
            &= \exp(\alpha \beta \log(x))
        \end{align*}
        as required.        
    \end{proof}
    
    \begin{corollary}
        Let \(p\), \(q \in \ZZ\), then
        \begin{align*}
            x^p = \underbrace{x \cdots x}_{p \text{ times}}
        \end{align*}
        and
        \begin{align*}
            x^{\frac p q} = \sqrt[q]{x^p}
        \end{align*}
        so our new definition matches our normal definition of exponentiation on exponents in \(\QQ\).
    \end{corollary}
    
    \begin{corollary}
        We have, for \(x \in \QQ\), \(\exp(x) = e^x\) in the normal sense.
    \end{corollary}
    
    \begin{proposition}[Growth of exponentials, powers and logarithms]
        Let \(r \in \RR^+\). Then
        \begin{itemize}
            \item \(x^r e^{-x} \to 0\) as \(x \to \infty\);
            \item \(x^{-r} \log x \to 0\) as \(x \to \infty\);
            \item \(x^r \log x \to 0^+\) as \(x \to 0^+\).
        \end{itemize}
    \end{proposition}
    
    \begin{proof}
        \begin{itemize}
            \item Let \(x > 1\), then
            \begin{align*}
                e^x = \sum_{k\geq 0} \frac{x^k}{k!} > \frac{x^n}{n!}
            \end{align*}
            for any \(n\). Choose \(n-r \geq 1\), then \(0 \leq x^r e^{-x} \leq \frac{n!}{x^{n-r}} \leq \frac{n!}{x}\), so \(x^r e^{-x} \to 0\) by sandwich.         
            \item We have \(t^{-1} \leq t^{\eps-1}\) for \(t \geq 1\) and any \(\eps > 0\). Pick \(\eps \in (0, r)\), so
            \begin{align*}
                0 \leq x^{-r} \log x &= x^-r \int_1^x \frac 1 t \D t  
                \\
                &\leq x^{-r} \int_1^x t^{\eps-1} \D t \\
                &\leq \frac 1 \eps x^{\eps-r}
            \end{align*}
            which goes to 0 as \(x \to \infty\), giving the result.
            \item By putting \(x = e^{-t}\) and \(y=rt\), we have
            \begin{align*}
                \lim_{x \to 0^+} x^r \log x = -\lim_{t \to \infty} -t e^{-tr} = -\lim_{y \to \infty} \frac 1 r ye^{-y} = 0
            \end{align*}
        \end{itemize}
    \end{proof}

    \subsection{Trigonometric functions}

    \begin{definition}[Sine and cosine]
        We define the functions sine and cosine as follows:
        \begin{align*}
            \cos: \CC &\to \CC \\
            z &\mapsto \frac 1 2 \left(e^{iz}+e^{-iz}\right) = \sum_{k=0}^\infty \frac{(-1)^k}{(2k!)}z^{2k} \\[10pt]
            \sin: \CC &\to \CC \\
            z &\mapsto \frac 1 {2i} \left(e^{iz}-e^{-iz}\right) = \sum_{k=0}^\infty \frac{(-1)^k}{(2k+1)!}z^{2k+1}
        \end{align*}
    \end{definition}
    
    \begin{lemma}
        We have, for \(w\), \(z \in \CC\),
        \begin{itemize}
            \item \(\cos(0) = 1\), \(\sin(0) = 0\);
            \item \(\cos'(z) = \sin z\), \(\sin'(z) = \cos(z)\);
            \item \(\sin(z+w) = \sin z \cos w + \cos z \sin w\);
            \item \(\cos(z+w) = \cos z \cos w - \sin z \sin w\);
            \item \(\sin^2 z + \cos^2 z = 1\)
        \end{itemize}
    \end{lemma}
    
    \begin{proof}
        This can all be derived with relative ease. For example,
        \begin{align*}
            \cos z \cos w - \sin z \sin w &= \frac 1 4 \left(e^{iz}+e^{-iz}\right)\left(e^{iw}+e^{-iw}\right) + \frac 1 4 \left(e^{iz}-e^{-iz}\right)\left(e^{iw}-e^{-iw}\right) \\
            &= \frac{1}{2} \left(e^{iz}e^{iw}+e^{-iz}e^{-iw}\right) = \cos(z+w)
        \end{align*}
        Taking now \(z=-w\), we can obtain the last result.
    \end{proof}
    
    On the reals, the identity \(\sin^2 x + \cos^2 x = 1\) gives that \(|\sin x|\), \(|\cos x| \leq 1\).

    \begin{proposition}[Periodicity of sine and cosine]
        There exists a smallest positive value \(\omega\) such that \(\cos(\omega/2) = 0\), \(\sin(\omega/2)=1\), and
        \begin{itemize}
            \item \(\sin\) and \(\cos\) are periodic with period \(2\omega\);
            \item \(\sin(x+\omega) = -\sin x\), \(\cos(x+\omega)=-\cos x\);
            \item \(\sin(x+\omega/2) = \cos x\), \(\cos(x+\omega/2)=-\sin x\).
        \end{itemize}
    \end{proposition}
    
    \begin{proof}
        Once we have shown the very first part of the statement, the rest of the results follow by the addition theorems. Let us look at \(x \in (0, 2)\), and note that
        \begin{align*}
            \cos(x) = -\sin(x) = -\left[\left(x-\frac{x^3}{3!}\right)+\left(\frac{x^5}{5!}-\frac{x^7}{7!}\right)+\dots\right]
        \end{align*}
        since we can rearrange and combine terms by the absolute convergence of the series. Examining each of these parirs of terms, they are all positive, so we conclude \(\cos\) is decreasing on this interval (and \(\sin\) is positive).

        It follows \(\cos\) has at most one root in this interval by Rolle's theorem. To see the existence of the root, observe
        \begin{align*}
            \cos(\sqrt 2) &= \underbrace{\left(1 - \frac{(\sqrt 2)^2}{2!}\right)}_{>0} + \underbrace{\left(\frac{(\sqrt 2)^4}{4!} - \frac{(\sqrt 2)^6}{6!}\right)}_{>0} + \dots \\
            \cos(\sqrt 3) &= \underbrace{1 - \frac{(\sqrt 3)^2}{2!} + \frac{(\sqrt 3)^4}{4!}}_{-\frac 1 8} - \underbrace{\left(\frac{(\sqrt 3)^6}{6!}-\frac{(\sqrt 3)^8}{8!}\right)}_{>0} - \dots
        \end{align*}
        so that \(\cos(\sqrt 2) > 0\) and \(\cos(\sqrt 3) < 0\). Hence, by the intermediate value theorem, we have \(\omega \in (\sqrt 2, \sqrt 3) \subset (0, 2)\) with \(\cos(\omega/2) = 0\). Hence, \(\sin(\omega/2) = \pm 1\), but we know \(\sin\) is positive, so \(\sin(\omega/2)=1\).
    \end{proof}
    
    \begin{corollary}
        We have that
        \begin{itemize}
            \item For \(x \in \RR\), the function \(e^{ix}\) is periodic with period \(2\omega\);
            \item \(e^{i\omega} = -1\);
            \item \(e^{i\omega/2} = i\).
        \end{itemize}
    \end{corollary}
    
    \begin{lemma}
        \(2\omega\) is the perimeter of the unit circle, \(S^1 = \{|z| = 1: z \in \CC\}\).
    \end{lemma}
    
    \begin{proof}
        We claim that
        \begin{align*}
            \gamma: [0, 2\omega) &\to S^1 \\
            t &\mapsto e^{it}
        \end{align*}
        is a bijection. We show
        \begin{itemize}
            \item \emph{\(\gamma\) is surjective.} Take \(z \in \S^1\) with \(z=a+ib\), \(a\), \(b \in \RR\), and \(a^2+b^2=1\). Hence \(|a|\), \(|b| \leq 1\), so, since \(c: [0, \omega] \to [-1, 1]\) is continuous and strictly monotone, we have \(t \in [0, \omega]\) with \(c(t) = a\).
            
            Furthermore, since \(s: [0, \omega] \to [-1, 1]\) is non-negative, we have \(\sin(t) = \sqrt{1-a^2}\). Hence, if \(b=\sqrt{1-a^2}\), we get the result as desired. On the other hand, if \(b=-\sqrt{1-a^2}\), we can make \(b=-\sin(t)=\sin(2\omega-t)\), and use that \(\cos(t)=\cos(2\omega-t)\). This gives the result.

            \item \emph{\(\gamma\) is injective.} Since \(e^{it}\) has least positive period \(2\omega\) by definition, it is injective on \([0, 2\omega)\).
        \end{itemize}

        Now, the perimeter of \(S^1\) is the length of \(\gamma\), i.e.
        \begin{align*}
            \int_0^{2\omega} |\gamma'(t)| \D t &= \int_0^{2\omega} |ie^{it}| \D t \\
            &= \int_0^{2\omega} \D t = 2\omega
        \end{align*}
    \end{proof}
    
    \begin{definition}[Hyperbolic sine and cosine]
        We define the hyperbolic sine and cosine functions as follows:
        \begin{align*}
            \cosh: \CC &\to \CC \\
            z &\mapsto \frac 1 2 \left(e^z+e^{-z}\right) = \cos(iz) \\[10pt]
            \sinh: \CC &\to \CC \\
            z &\mapsto \frac 1 2 \left(e^z-e^{-z}\right) = -i\sin(iz)
        \end{align*}
    \end{definition}
    
    \begin{lemma}
        We have, for \(z \in \CC\),
        \begin{itemize}
            \item \(\cosh'z=\sinh z\);
            \item \(\sinh'z=\cosh z\);
            \item \(\cosh^2z-\sinh^2z=1\).
        \end{itemize}
    \end{lemma}
    
    
    
    
    
    
    
    
    


\end{document}
